%Version 3 October 2023
% See section 11 of the User Manual for version history
%
%%%%%%%%%%%%%%%%%%%%%%%%%%%%%%%%%%%%%%%%%%%%%%%%%%%%%%%%%%%%%%%%%%%%%%
%%                                                                 %%
%% Please do not use \input{...} to include other tex files.       %%
%% Submit your LaTeX manuscript as one .tex document.              %%
%%                                                                 %%
%% All additional figures and files should be attached             %%
%% separately and not embedded in the \TeX\ document itself.       %%
%%                                                                 %%
%%%%%%%%%%%%%%%%%%%%%%%%%%%%%%%%%%%%%%%%%%%%%%%%%%%%%%%%%%%%%%%%%%%%%

%%\documentclass[referee,sn-basic]{sn-jnl}% referee option is meant for double line spacing

%%=======================================================%%
%% to print line numbers in the margin use lineno option %%
%%=======================================================%%

%%\documentclass[lineno,sn-basic]{sn-jnl}% Basic Springer Nature Reference Style/Chemistry Reference Style

%%======================================================%%
%% to compile with pdflatex/xelatex use pdflatex option %%
%%======================================================%%

%%\documentclass[pdflatex,sn-basic]{sn-jnl}% Basic Springer Nature Reference Style/Chemistry Reference Style


%%Note: the following reference styles support Namedate and Numbered referencing. By default the style follows the most common style. To switch between the options you can add or remove “Numbered” in the optional parenthesis. 
%%The option is available for: sn-basic.bst, sn-vancouver.bst, sn-chicago.bst%  
 
\documentclass[sn-nature,lineno]{sn-jnl}% Style for submissions to Nature Portfolio journals
%%\documentclass[sn-basic]{sn-jnl}% Basic Springer Nature Reference Style/Chemistry Reference Style
%%\documentclass[sn-mathphys-num]{sn-jnl}% Math and Physical Sciences Numbered Reference Style 
%%\documentclass[sn-mathphys-ay]{sn-jnl}% Math and Physical Sciences Author Year Reference Style
%%\documentclass[sn-aps]{sn-jnl}% American Physical Society (APS) Reference Style
%%\documentclass[sn-vancouver,Numbered]{sn-jnl}% Vancouver Reference Style
%%\documentclass[sn-apa]{sn-jnl}% APA Reference Style 
%%\documentclass[sn-chicago]{sn-jnl}% Chicago-based Humanities Reference Style

%%%% Standard Packages
%%<additional latex packages if required can be included here>

\usepackage{graphicx}%
\usepackage{multirow}%
\usepackage{amsmath,amssymb,amsfonts}%
\usepackage{amsthm}%
\usepackage{mathrsfs}%
\usepackage[title]{appendix}%
\usepackage{xcolor}%
\usepackage{textcomp}%
\usepackage{manyfoot}%
\usepackage{booktabs}%
\usepackage{algorithm}%
\usepackage{algorithmicx}%
\usepackage{algpseudocode}%
\usepackage{listings}%
\usepackage{subcaption}
\usepackage{comment}
\usepackage{tikz}
\usepackage{luacode}
\usepackage{pdfpages}
\usepackage{hyperref}
%%%%

%%%%%=============================================================================%%%%
%%%%  Remarks: This template is provided to aid authors with the preparation
%%%%  of original research articles intended for submission to journals published 
%%%%  by Springer Nature. The guidance has been prepared in partnership with 
%%%%  production teams to conform to Springer Nature technical requirements. 
%%%%  Editorial and presentation requirements differ among journal portfolios and 
%%%%  research disciplines. You may find sections in this template are irrelevant 
%%%%  to your work and are empowered to omit any such section if allowed by the 
%%%%  journal you intend to submit to. The submission guidelines and policies 
%%%%  of the journal take precedence. A detailed User Manual is available in the 
%%%%  template package for technical guidance.
%%%%%=============================================================================%%%%

%% as per the requirement new theorem styles can be included as shown below
\theoremstyle{thmstyleone}%
\newtheorem{theorem}{Theorem}%  meant for continuous numbers
%%\newtheorem{theorem}{Theorem}[section]% meant for sectionwise numbers
%% optional argument [theorem] produces theorem numbering sequence instead of independent numbers for Proposition
\newtheorem{proposition}[theorem]{Proposition}% 
%%\newtheorem{proposition}{Proposition}% to get separate numbers for theorem and proposition etc.

\theoremstyle{thmstyletwo}%
\newtheorem{example}{Example}%
\newtheorem{remark}{Remark}%

\theoremstyle{thmstylethree}%
\newtheorem{definition}{Definition}%

%% Import variables

\begin{luacode*}
    -- Define dictionary of variables below.
        dictionary = {
		var__OTG_GWAS_SE = "0.30",
		var__PIGEAN_GWAS_SE = "0.53",
		var__OTG_PIGEAN_COMBINED_SE = "0.61",
		var__OTG_PIGEAN_COMBINED_Z_SCORE = "2.03",
		var__OTG_PIGEAN_COMBINED_P_VALUE = "0.0212",
		var__CCAT_EXACT_GENEBASS_ALL_GENEBASS_ALL_AREAS = "3.13",
		var__CCAT_CI_LOWER_GENEBASS_ALL_GENEBASS_ALL_AREAS = "1.95",
		var__CCAT_CI_UPPER_GENEBASS_ALL_GENEBASS_ALL_AREAS = "5.02",
		var__BASELINE_CCAT_EXACT_GENEBASS_ALL_GENEBASS_ALL_AREAS = "1.06",
		var__BASELINE_CCAT_CI_LOWER_GENEBASS_ALL_GENEBASS_ALL_AREAS = "0.93",
		var__BASELINE_CCAT_CI_UPPER_GENEBASS_ALL_GENEBASS_ALL_AREAS = "1.22",
		var__CCAT_EXACT_ONCOLOGY_ALL_INTOGEN_ALL_AREAS = "nan",
		var__CCAT_CI_LOWER_ONCOLOGY_ALL_INTOGEN_ALL_AREAS = "nan",
		var__CCAT_CI_UPPER_ONCOLOGY_ALL_INTOGEN_ALL_AREAS = "nan",
		var__BASELINE_CCAT_EXACT_CONGENITAL_ALL_INTOGEN_ALL_AREAS = "nan",
		var__BASELINE_CCAT_CI_LOWER_CONGENITAL_ALL_INTOGEN_ALL_AREAS = "nan",
		var__BASELINE_CCAT_CI_UPPER_CONGENITAL_ALL_INTOGEN_ALL_AREAS = "nan",
		var__CCAT_EXACT_OMIM_ALL_ALL_AREAS = "3.74",
		var__CCAT_CI_LOWER_OMIM_ALL_ALL_AREAS = "3.06",
		var__CCAT_CI_UPPER_OMIM_ALL_ALL_AREAS = "4.57",
		var__CCAT_EXACT_OMIM_PIGEAN_FILTERED_ALL_AREAS = "4.52",
		var__CCAT_CI_LOWER_OMIM_PIGEAN_FILTERED_ALL_AREAS = "2.14",
		var__CCAT_CI_UPPER_OMIM_PIGEAN_FILTERED_ALL_AREAS = "9.53",
		var__BASELINE_CCAT_EXACT_OMIM_ALL_ALL_AREAS = "1.27",
		var__BASELINE_CCAT_CI_LOWER_OMIM_ALL_ALL_AREAS = "1.16",
		var__BASELINE_CCAT_CI_UPPER_OMIM_ALL_ALL_AREAS = "1.39",
		var__BASELINE_CCAT_EXACT_OMIM_PIGEAN_FILTERED_ALL_AREAS = "1.58",
		var__BASELINE_CCAT_CI_LOWER_OMIM_PIGEAN_FILTERED_ALL_AREAS = "1.28",
		var__BASELINE_CCAT_CI_UPPER_OMIM_PIGEAN_FILTERED_ALL_AREAS = "1.95",
		var__CCAT_EXACT_OTG_FINNGEN_ALL_AREAS = "3.02",
		var__CCAT_CI_LOWER_OTG_FINNGEN_ALL_AREAS = "1.72",
		var__CCAT_CI_UPPER_OTG_FINNGEN_ALL_AREAS = "5.29",
		var__CCAT_EXACT_OTG_GWAS_CATALOG_ALL_AREAS = "2.11",
		var__CCAT_CI_LOWER_OTG_GWAS_CATALOG_ALL_AREAS = "1.60",
		var__CCAT_CI_UPPER_OTG_GWAS_CATALOG_ALL_AREAS = "2.77",
		var__CCAT_EXACT_OTG_NEALE_UKBB_ALL_AREAS = "4.07",
		var__CCAT_CI_LOWER_OTG_NEALE_UKBB_ALL_AREAS = "2.03",
		var__CCAT_CI_UPPER_OTG_NEALE_UKBB_ALL_AREAS = "8.14",
		var__CCAT_EXACT_OTG_ALL_ALL_AREAS = "1.98",
		var__CCAT_CI_LOWER_OTG_ALL_ALL_AREAS = "1.70",
		var__CCAT_CI_UPPER_OTG_ALL_ALL_AREAS = "2.31",
		var__CCAT_EXACT_OTG_SAIGE_ALL_AREAS = "2.71",
		var__CCAT_CI_LOWER_OTG_SAIGE_ALL_AREAS = "0.96",
		var__CCAT_CI_UPPER_OTG_SAIGE_ALL_AREAS = "7.61",
		var__BASELINE_CCAT_EXACT_OTG_FINNGEN_ALL_AREAS = "0.78",
		var__BASELINE_CCAT_CI_LOWER_OTG_FINNGEN_ALL_AREAS = "0.68",
		var__BASELINE_CCAT_CI_UPPER_OTG_FINNGEN_ALL_AREAS = "0.89",
		var__BASELINE_CCAT_EXACT_OTG_GWAS_CATALOG_ALL_AREAS = "0.74",
		var__BASELINE_CCAT_CI_LOWER_OTG_GWAS_CATALOG_ALL_AREAS = "0.67",
		var__BASELINE_CCAT_CI_UPPER_OTG_GWAS_CATALOG_ALL_AREAS = "0.81",
		var__BASELINE_CCAT_EXACT_OTG_NEALE_UKBB_ALL_AREAS = "1.01",
		var__BASELINE_CCAT_CI_LOWER_OTG_NEALE_UKBB_ALL_AREAS = "0.81",
		var__BASELINE_CCAT_CI_UPPER_OTG_NEALE_UKBB_ALL_AREAS = "1.25",
		var__BASELINE_CCAT_EXACT_OTG_ALL_ALL_AREAS = "0.68",
		var__BASELINE_CCAT_CI_LOWER_OTG_ALL_ALL_AREAS = "0.62",
		var__BASELINE_CCAT_CI_UPPER_OTG_ALL_ALL_AREAS = "0.76",
		var__BASELINE_CCAT_EXACT_OTG_SAIGE_ALL_AREAS = "1.34",
		var__BASELINE_CCAT_CI_LOWER_OTG_SAIGE_ALL_AREAS = "1.09",
		var__BASELINE_CCAT_CI_UPPER_OTG_SAIGE_ALL_AREAS = "1.65",
		var__CCAT_EXACT_PICCOLO_ALL_PICCOLO_ALL_AREAS = "2.11",
		var__CCAT_CI_LOWER_PICCOLO_ALL_PICCOLO_ALL_AREAS = "1.53",
		var__CCAT_CI_UPPER_PICCOLO_ALL_PICCOLO_ALL_AREAS = "2.91",
		var__BASELINE_CCAT_EXACT_PICCOLO_ALL_PICCOLO_ALL_AREAS = "0.70",
		var__BASELINE_CCAT_CI_LOWER_PICCOLO_ALL_PICCOLO_ALL_AREAS = "0.64",
		var__BASELINE_CCAT_CI_UPPER_PICCOLO_ALL_PICCOLO_ALL_AREAS = "0.77",
		var__CCAT_EXACT_PIGEAN_ALL_ALL_AREAS = "2.93",
		var__CCAT_CI_LOWER_PIGEAN_ALL_ALL_AREAS = "2.45",
		var__CCAT_CI_UPPER_PIGEAN_ALL_ALL_AREAS = "3.51",
		var__CCAT_EXACT_PIGEAN_GWAS_CATALOG_ALL_AREAS = "3.35",
		var__CCAT_CI_LOWER_PIGEAN_GWAS_CATALOG_ALL_AREAS = "2.46",
		var__CCAT_CI_UPPER_PIGEAN_GWAS_CATALOG_ALL_AREAS = "4.55",
		var__CCAT_EXACT_PIGEAN_OMIM_ALL_AREAS = "3.22",
		var__CCAT_CI_LOWER_PIGEAN_OMIM_ALL_AREAS = "2.07",
		var__CCAT_CI_UPPER_PIGEAN_OMIM_ALL_AREAS = "5.01",
		var__CCAT_EXACT_PIGEAN_ORPHANET_ALL_AREAS = "2.76",
		var__CCAT_CI_LOWER_PIGEAN_ORPHANET_ALL_AREAS = "2.23",
		var__CCAT_CI_UPPER_PIGEAN_ORPHANET_ALL_AREAS = "3.41",
		var__CCAT_EXACT_PIGEAN_A2F_PORTAL_ALL_AREAS = "3.50",
		var__CCAT_CI_LOWER_PIGEAN_A2F_PORTAL_ALL_AREAS = "2.32",
		var__CCAT_CI_UPPER_PIGEAN_A2F_PORTAL_ALL_AREAS = "5.29",
		var__BASELINE_CCAT_EXACT_PIGEAN_ALL_ALL_AREAS = "0.81",
		var__BASELINE_CCAT_CI_LOWER_PIGEAN_ALL_ALL_AREAS = "0.73",
		var__BASELINE_CCAT_CI_UPPER_PIGEAN_ALL_ALL_AREAS = "0.90",
		var__BASELINE_CCAT_EXACT_PIGEAN_GWAS_CATALOG_ALL_AREAS = "0.77",
		var__BASELINE_CCAT_CI_LOWER_PIGEAN_GWAS_CATALOG_ALL_AREAS = "0.70",
		var__BASELINE_CCAT_CI_UPPER_PIGEAN_GWAS_CATALOG_ALL_AREAS = "0.85",
		var__BASELINE_CCAT_EXACT_PIGEAN_OMIM_ALL_AREAS = "1.48",
		var__BASELINE_CCAT_CI_LOWER_PIGEAN_OMIM_ALL_AREAS = "1.22",
		var__BASELINE_CCAT_CI_UPPER_PIGEAN_OMIM_ALL_AREAS = "1.79",
		var__BASELINE_CCAT_EXACT_PIGEAN_ORPHANET_ALL_AREAS = "1.06",
		var__BASELINE_CCAT_CI_LOWER_PIGEAN_ORPHANET_ALL_AREAS = "0.97",
		var__BASELINE_CCAT_CI_UPPER_PIGEAN_ORPHANET_ALL_AREAS = "1.17",
		var__BASELINE_CCAT_EXACT_PIGEAN_A2F_PORTAL_ALL_AREAS = "0.86",
		var__BASELINE_CCAT_CI_LOWER_PIGEAN_A2F_PORTAL_ALL_AREAS = "0.78",
		var__BASELINE_CCAT_CI_UPPER_PIGEAN_A2F_PORTAL_ALL_AREAS = "0.94",
		var__PIGEAN_COMBINED_REGRESSION_R2 = "0.85",
		var__PIGEAN_COMBINED_REGRESSION_P_VALUE = "0.0002",
		var__PIGEAN_COMBINED_REGRESSION_COEF = "0.13",
		var__NUM_PORTAL_PIGEAN_TRAITS = "1,011",
		var__NUM_ORPHANET_PIGEAN_TRAITS = "1,475",
		var__NUM_GCAT_PIGEAN_TRAITS = "4,002",
		var__NUM_PIGEAN_TRAITS = "6,488",
		var__NUM_OMIM_PIGEAN_TRAITS = "100",
		var__PIGEAN_DIRECT_REGRESSION_R2 = "0.94",
		var__PIGEAN_DIRECT_REGRESSION_P_VALUE = "7.5142e-13",
		var__PIGEAN_DIRECT_REGRESSION_COEF = "0.29",
		var__PIGEAN_DIRECT_MERGED_COMMON_REGRESSION_R2 = "0.65",
		var__PIGEAN_DIRECT_MERGED_COMMON_REGRESSION_P_VALUE = "9.8037e-06",
		var__PIGEAN_DIRECT_MERGED_COMMON_REGRESSION_COEF = "0.74",
		var__PIGEAN_DIRECT_MERGED_RARE_REGRESSION_R2 = "0.22",
		var__PIGEAN_DIRECT_MERGED_RARE_REGRESSION_P_VALUE = "0.0284",
		var__PIGEAN_DIRECT_MERGED_RARE_REGRESSION_COEF = "0.13",
		var__PIGEAN_RIDGE_REGRESSION_R2 = "0.60",
		var__PIGEAN_RIDGE_REGRESSION_R2_ADJ = "0.58",
		var__PIGEAN_INDIRECT_REGRESSION_R2 = "0.77",
		var__PIGEAN_INDIRECT_REGRESSION_P_VALUE = "2.0524e-07",
		var__PIGEAN_INDIRECT_REGRESSION_COEF = "0.13",
		var__PIGEAN_INDIRECT_MERGED_COMMON_REGRESSION_R2 = "0.94",
		var__PIGEAN_INDIRECT_MERGED_COMMON_REGRESSION_P_VALUE = "1.2283e-10",
		var__PIGEAN_INDIRECT_MERGED_COMMON_REGRESSION_COEF = "0.32",
		var__PIGEAN_INDIRECT_MERGED_RARE_REGRESSION_R2 = "0.68",
		var__PIGEAN_INDIRECT_MERGED_RARE_REGRESSION_P_VALUE = "5.0136e-05",
		var__PIGEAN_INDIRECT_MERGED_RARE_REGRESSION_COEF = "0.35",
		var__NUM_UNIQUE_DRUG_PROGRAMS = "29,460",
		var__NUM_UNIQUE_DRUG_TARGETS = "2,517",
		var__NUM_UNIQUE_INDICATIONS = "1,013",
		var__NUM_UNIQUE_GCAT_MAPPED_TRAITS = "986",
		var__NUM_UNIQUE_ORPHANET_MAPPED_TRAITS = "762",
		var__NUM_UNIQUE_PORTAL_TO_MESH_MAPPED_TRAITS = "397",
		var__NUM_UNIQUE_MAPPED_TRAITS = "1,744",
		var__NUM_UNIQUE_MINIKEL_ASSOCIATIONS = "184,839",

    }
    
    function lookup(key)
        if dictionary[key] then
            tex.print(dictionary[key])
        else
            tex.print("XXXX")
        end
    end
\end{luacode*}


%% New commands
% Probablity Dist function
\newcommand{\prob}{\text{Pr}}
% Variable lookup function
\newcommand{\getvar}[1]{\directlua{lookup("#1")}}

\raggedbottom
%%\unnumbered% uncomment this for unnumbered level heads

\begin{document}

\title[Article Title]{Beyond the Locus: Modeling Indirect Genetic Support to Expand the Reach of Human Genetics}

%Dual embeddings of 6,915 traits and 18,125 genes learned from direct and indirect human genetic support
%A dense statistical map of direct and indirect genetic support between 18,125 genes and 6,915 rare and common human traits


%%=============================================================%%
%% GivenName	-> \fnm{Joergen W.}
%% Particle	-> \spfx{van der} -> surname prefix
%% FamilyName	-> \sur{Ploeg}
%% Suffix	-> \sfx{IV}
%% \author*[1,2]{\fnm{Joergen W.} \spfx{van der} \sur{Ploeg} 
%%  \sfx{IV}}\email{iauthor@gmail.com}
%%=============================================================%%

\author*[1,2]{\fnm{First} \sur{Author}}\email{iauthor@gmail.com}

\author[2,3]{\fnm{Second} \sur{Author}}\email{iiauthor@gmail.com}
\equalcont{These authors contributed equally to this work.}

\author[1,2]{\fnm{Third} \sur{Author}}\email{iiiauthor@gmail.com}
\equalcont{These authors contributed equally to this work.}

\affil*[1]{\orgdiv{Department}, \orgname{Organization}, \orgaddress{\street{Street}, \city{City}, \postcode{100190}, \state{State}, \country{Country}}}

\affil[2]{\orgdiv{Department}, \orgname{Organization}, \orgaddress{\street{Street}, \city{City}, \postcode{10587}, \state{State}, \country{Country}}}

\affil[3]{\orgdiv{Department}, \orgname{Organization}, \orgaddress{\street{Street}, \city{City}, \postcode{610101}, \state{State}, \country{Country}}}

%%==================================%%
%% Sample for unstructured abstract %%
%%==================================%%

\abstract{
Human genetic evidence is known to increase the likelihood of clinical success in drug development \cite{ochoa_human_2022,rusina_genetic_2023}. Yet fewer than 5\% of active drug programs are supported by germline genetic data \cite{minikel_refining_2024}, revealing a critical gap in how genetic knowledge is incorporated into therapeutic pipelines. To help bridge this gap, we introduce the concept of indirect genetic support, which captures evidence derived from associations with biologically related genes linked through shared annotations. This expands on direct genetic support by probabilistically integrating both genetic and non-genetic information into a unified model aimed at improving mechanistic understanding and clinical outcomes. We present PIGEAN, a Bayesian framework that quantifies direct and indirect support using GWAS summary statistics, exome-wide gene-level associations, and rare disease gene sets. Applied across 6,000 complex traits, 2,000 rare diseases, 600,000 gene annotations, and 29,000 drug programs, PIGEAN constructs a weighted graph that serves as a genetically grounded backbone for biomedical discovery. Our analysis shows that drug targets prioritized via this graph achieve clinical success rates comparable to those found in curated drug target databases, while also revealing over 1,500 additional viable targets through indirect evidence alone. By organizing gene prioritization along the interpretable axes of direct and indirect support, PIGEAN offers a principled alternative to existing ad hoc methods. When paired with AI agents, this framework enables scalable, mechanistically informed exploration of the gene–phenotype landscape, probablistically connecting every gene to every trait using the structure of current biological knowledge.
}

%NOTES FOR ABSTRACT
%Novel inferences of course
%But beyond that, the presence of known and novel all in one place enables analyses (RS, gene function queries, gene lists for disease, gene lists for collection of phenotypes ) that could not be done previously
%In particular, shows shared basis of rare and common diseases, which increases further when considered at pathway level and not just genes


%%% Old abstract
% \abstract{Human genetic support increases the likelihood of therapeutic success and helps prioritize genes for further biological study. Here, we introduce the idea of indirect genetic support, which we distinguish from the traditional concept of direct support in that genetic associations for other genes support a target gene through shared gene annotations. We derive a mathemtical model, PIGEAN, that probabilistically estimates indirect genetic support from GWAS summary statisics, exome gene-level association statistics, and/or rare disease gene lists; we show that this model generalizes the idea of "bottom-up" or "convergence" based methods for gene prioritization and enables its combination with estimates of direct genetic support. We apply this model across \getvar{NUM_COMPLEX_TRAITS} complex traits, XXX rare diseases, and XXX gene annotations, assigning X\% of genes support for at least one trait. [Result from analysis of gene annotations]. [Result from analysis of how many GWAS loci have known mechanisms]. [Result from joint vs. marginal estimates and mechanisms]. [Result from drug target analysis]. These estimates of indirect support significantly expand the number of genes with links to human traits and enable probabilistic reasoning to replace ad hoc gene prioritization approaches that combine genetic associations with prior knowledge about gene function.
% }

%%================================%%
%% Sample for structured abstract %%
%%================================%%

% \abstract{\textbf{Purpose:} The abstract serves both as a general introduction to the topic and as a brief, non-technical summary of the main results and their implications. The abstract must not include subheadings (unless expressly permitted in the journal's Instructions to Authors), equations or citations. As a guide the abstract should not exceed 200 words. Most journals do not set a hard limit however authors are advised to check the author instructions for the journal they are submitting to.
% 
% \textbf{Methods:} The abstract serves both as a general introduction to the topic and as a brief, non-technical summary of the main results and their implications. The abstract must not include subheadings (unless expressly permitted in the journal's Instructions to Authors), equations or citations. As a guide the abstract should not exceed 200 words. Most journals do not set a hard limit however authors are advised to check the author instructions for the journal they are submitting to.
% 
% \textbf{Results:} The abstract serves both as a general introduction to the topic and as a brief, non-technical summary of the main results and their implications. The abstract must not include subheadings (unless expressly permitted in the journal's Instructions to Authors), equations or citations. As a guide the abstract should not exceed 200 words. Most journals do not set a hard limit however authors are advised to check the author instructions for the journal they are submitting to.
% 
% \textbf{Conclusion:} The abstract serves both as a general introduction to the topic and as a brief, non-technical summary of the main results and their implications. The abstract must not include subheadings (unless expressly permitted in the journal's Instructions to Authors), equations or citations. As a guide the abstract should not exceed 200 words. Most journals do not set a hard limit however authors are advised to check the author instructions for the journal they are submitting to.}

\keywords{keyword1, Keyword2, Keyword3, Keyword4}

%%\pacs[JEL Classification]{D8, H51}

%%\pacs[MSC Classification]{35A01, 65L10, 65L12, 65L20, 65L70}

\maketitle

\newpage
\tableofcontents
\newpage

\section{Main}\label{sec:introduction}

% Broad introductory paragraph
Most human traits are influenced by many genes, and most human genes influence many traits. Mapping this complex web of gene–trait relationships is essential for uncovering disease biology \cite{cohen_sequence_2006}, understanding clinical heterogeneity \cite{udler_type_2018}, accurately diagnosing patients \cite{yang_clinical_2013}, and designing effective therapies \cite{plenge_validating_2013}. Among the most widely used and powerful approaches to to identify gene-trait relationships are genetic association studies. For common traits, genome-wide association studies (GWAS) have identified tens of thousands of trait-associated loci across thousands of phenotypes \cite{abdellaoui_15_2023}, while rare disease studies have linked thousands of Mendelian conditions to specific causal genes \cite{chung_meta-analysis_2023}. A genetic association provides a gene with ``genetic support" for a human trait -- establishing the human-relevance and causal effect of its perturbation.

The number of trait-associated genes identified by genetic association studies is, however, but a small fraction of the total number of trait-relevant genes -- that is, those that can be perturbed in humans to yield an effect on a trait of interest\cite{finan_druggable_2017}. Many trait-relevant genes escape detection due to limited statistical power \cite{wang_statistical_2019}, constraining evolutionary pressure on genetic variation \cite{gazal_linkage_2017}, or the difficulty of assigning noncoding GWAS signals to their effector genes \cite{schipper_demystifying_2022,costanzo_realizing_2025}. Yet the full set of trait-relevant genes is of most interest for drug development and experimental study\cite{jhamb_pathway_2019,macnamara_network_2020}. In part due to this incompleteness, most translational research does not use human genetics to prioritize genes\cite{wehling_assessing_2009,dornbos_evaluating_2022}, even with the  increasing recognition that drug targets with genetic support are substantially more likely to succeed clinically\cite{plenge_validating_2013,minikel_refining_2024} 

% Related Work and More specific stage setting to elicit the the contributions paragraph
By contrast, non-human-genetic approaches are more widely used to identify trait-relevant genes but come with critical tradeoffs. Experimental studies in animal \cite{seok_genomic_2013,widjaja_inhibiting_2019} and cell-based \cite{ben-david_genetic_2018,canon_clinical_2019} systems can establish causal links between genes and phenotypes, but their relevance to humans is often uncertain. Human observational studies, such as gene expression profiling \cite{wu_integrating_2022,rodriguez_machine_2021} or epidemiological associations \cite{barter_cholesteryl_2011,bhatt_cardiovascular_2019}, provide human-relevant context but lack causal directionality. The literature and biomedical databases are rife with gene-trait relationships derived from these or other non-genetic studies, and they are widely used as a starting point for drug discovery or experimental study\cite{swinney_how_2011,dornbos_evaluating_2022}. Recent studies demonstrate how these sources can be rapidly and accurately queried with large language models (LLMs) \cite{toufiq_harnessing_2023,hu_evaluation_2025}. However, it is well documented that prioritizing genes using these evidence sources is prone to bias, incompleteness, or circularity\cite{brito_removing_2011,gillis_assessing_2013,gaudet_removing_2017,dunham_human_2018,haynes_gene_2018,stoeger_large_2018,richardson_meta_2024,tan_caution_2025}.

Ideally, a map of gene-trait relationships would combine the complementary benefits of human genetic and non-human-genetic evidence sources. It is commonplace to evaluate genes individually through manual expert synthesis of genetic data and the published literature\cite{mountjoy_open_2022,karjalainen_genome-wide_2024}, but a quantiative map across the genome and phenome requires scalabale statistical methods. Such methods have been developed both in the context of GWAS effector gene identification, where gene ``functional annotations'' enriched for genetic associations are used to prioritize genes nearby associated SNPs\cite{pers_biological_2015,weeks_leveraging_2023}, and in the context of rare disease diagnostics, where literature mining \cite{birgmeier_amelie_2020} and interaction networks\cite{smedley_next-generation_2015} are used to prioritize genes that carry variants observed in a patient. These approaches, however, rely on heuristics (such as harmonic averaging \cite{buniello_open_2025} or regressions using noisy or incomplete training data\cite{weeks_leveraging_2023,schipper_prioritizing_2025}) and often lack a biological model, which limits their portability across traits or application to settings other than that for which they were originally designed. As a result, the currently charted map of trait-gene relationships across rare and common diseases remains sparse, inaccurate, and inconsistent across traits.

% Contribution explanation
Here, we combine human genetic association data for rare and common traits and non-human-genetic data encoded in 600,000 gene annotations to construct a complete probabilistic evidence map relating 18,125 genes to 6,915 traits. Each gene-trait relationship is a composite of two statistically calibrated values: a traditional measure of  ``direct genetic support'' (the strength of gene-trait association from human genetic data) and a new  measure of ``indirect genetic support'' (the extent to which the gene shares functional annotations with other trait-associated genes). Indirect genetic support is measured on the same scale as direct genetic support, estimates a measure of gene-trait relevance distinct from direct gene-trait association, and enables us to use unbiased human genetic associations to calibrate noisy and heterogeneous non-human-genetic data sources and learn which are trustworthy for each trait. This map when thresholded is orders of magnitude denser than manually curated knowledge graphs and, when unthresholded, can be viewed as dual numerical embeddings of genes and traits learned from combined genetic and non-genetic data sources. We illustrate the utility of this resource to aid in the interpretation of gene function and trait mechanistic heterogeneity, predict drug target relative success more accurately and comprehensively than previous resources, and provide a resource for agentic AI not capturable by pretraining on text-only corpora.

%This decomposition formalizes and generalizes previous heuristic or machine learning-based integration strategies \cite{smedley_next-generation_2015,weeks_leveraging_2023,schipper_prioritizing_2025}, offering a statistically grounded alternative.

%Unlike conventional biomedical knowledge graphs \cite{li_graph_2022}, the PGEG avoids lossy edge labeling \cite{himmelstein_systematic_2017,morris_scalable_2023}, heuristic integration \cite{ochoa_next-generation_2023}, and manual curation \cite{putman_monarch_2024}, while remaining directly usable for downstream applications.

\section{Results}
\subsection{Overview of methods}
\begin{figure}[h!]
    \centering
    \includegraphics[width=.85\linewidth]{figures/PIGEAN-Figure-1.pdf}
    \caption{Conceptual overview of the PIGEAN framework. (A) Demostration of top-down prioritization methods from GWAS, WES and/or Mendelian disease databases. (B) Demostration of bottom-up priortization methods from gene annotations. (C) Graphical Model of PIGEAN with a depiction of fusing direct and indirect genetic evidence. (D) Positioning each gene across the continum of direct and indirect support. (E) Leveraging this continuum to power downstream analysis pipelines which can be leveraged in agentic tooling (F) to yeild AI-driven analysis at scale.}
    \label{fig:pigean-concept}
\end{figure}
To map gene-trait relationships across the entire phenome and genome, we sought to infer a probability of relevance for any gene  to any trait. Doing so using available human genetic associations and non-human-genetic data posed several challenges. First, we needed to infer relationships across thousands of traits accurately, robustly, and efficiently. Second, the traits are heterogeneous in terms of genetic association data available (e.g. common variant associations for GWAS, Mendelian gene-disease validity assertions for rare diseases) and genetic architecture. Third, currently significant gene-trait associations cover only a small fraction of trait-relevant genes. Fourth, non-human-genetic datasets (i.e. curated pathway databases, results from animal model experiments, text-mined relationships, or gene expression enrichment analyses) are large in number and of varying quality, and it is unclear which datasets are useful for predicting trait-relevant genes. Finally, non-human-genetic datasets are biased toward well-studied genes and well-studied biological mechanisms.

To address these challenges, we developed a statistical model of a quantity we term a gene's ``indirect genetic support'', and a method termed PIGEAN (Priors Inferred from GEne ANnotations) for homogeneously inferring across common and rare traits. Indirect genetic support combines the unbiased rigor of genetic associations with the breadth and scope of gene annotations, which encode non-genetic-data either as gene sets (binary annotations) or gene quantitative vectors (continuous annotations) and can be constructed in a high-throughput manner from many biomedical data sources (\ref{sec:methods}). Indirect genetic support is an emergent property of a simple statistical model of observed genetic associations, in which gene annotations affect the probability that a gene is relevant to a trait, with effect sizes unknown and different for each annotation and each trait, and in which trait-relevant genes are more likely to produce genetic associations than are non-trait-relevant genes (\ref{sec:methods}).  This model statistically generalizes several widely used heuristic approaches for prioritizing genes from genetic association studies (\ref{sec:methods}), and it enables us to demonstrate mathematically that genetic associations not only directly implicate the genes for which they are observed but also indirectly implicate genes with shared annotations. PIGEAN applies Bayesian inference to this model to jointly infer gene annotation effect sizes and gene trait-relevance probabilities given a library of gene annotations and either GWAS summary statistics or rare disease gene lists, in so doing (a) using human genetic associations to select annotations with non-zero effects and calibrate their effect sizes and (b) using genetic annotations to identify additional trait-relevant genes beyond those with direct association evidence. Priors on the latent parameters enforce strict regularization that control biases towards well-studied genes or profligate annotations, as well as enforcing model sparsity and interpretability (\ref{sec:methods}).

We studied PIGEAN across a universe of 5,030 complex traits and 1,526 rare disease groups. For complex traits, we integrated full summary statistic files from XXX GWAS covering XXX traits, as well as curated association lists from XXX GWAS across XXX traits. We used an extension of a previously published approach \cite{dornbos_evaluating_2022} to calculate each gene's direct genetic support from variant-level associations (\ref{sec:methods}).  For rare diseases, we mapped curated disease–gene confidence levels from Orphanet to probabilities of disease relevance (\ref{sec:methods}).  

To evaluate the contribution of non-genetic evidence, we constructed multiple PIGEAN models, each defined by a specific set of geneset libraries available during inference. These libraries collectively included 144,627 gene annotations spanning diverse sources (\ref{Supplementary Table 2}): knockout mouse phenotypes (11,502 Mammalian Phenotype Ontology terms), membership in pathways and curated gene sets (29,003 from the Molecular Signatures Database), co-occurrence in PubMedCentral manuscript figures (19,391 sets from the Pathway Figure OCR project), co-occurrence with terms in PubMed abstracts (25,126 MeSH terms), network interaction partners (37,156 STRING neighbors), and single-cell gene expression patterns (41,840 experimental clusters curated as part of the PoPS method).  

The resulting map (Figure \ref{fig:full_map}), generated from the top-performing model identified via model selection (See \textit{Validation of Resource}), produced a complete numerical matrix of gene–trait relevance probabilities. This included 135,188 gene-trait relationships (21 per trait and 7 per gene) with a posterior probability (PP) above 90\% (i.e. false discovery rate < 0.1), 295,547 gene–trait relationships (45 per trait and 16 per gene) with PP>50\%, and 10.49M significant (\ref{sec:methods}) gene set–trait associations. Each gene-trait relationship is supported by provenance that includes the specific SNP associations or curated gene records  that produce direct support, and the specific gene set-trait relationsips that produce indirect support. The full map is available for download, for interactive browsing via a web interface, or through AI-driven query to filter and explore results in response to high-level questions.  

\subsection{Gene-trait relationships with direct or indirect genetic support}

%Some stat for each paragraph to describe the properties of the genes
%Highest combined: look at network centrality? Conserved?
%Common suggestive: how many near signals vs not? How many increased in gwas signals vs not?
%How constrained as a function of support? Direct vs indirect which is stronger? Condition on combined and then residual on indirect?
%Try pli and Loef
%Try to make the paragraphs each say something about what indirect is telling you about the gene, disease relevant genes vs disease associated genes, etc. then the examples make it concrete 

%Points in this section
%We have the super strong positive controls?
%Then we have the 80%-90%. Is this worth a new paragraph or just combining with the 99%. These are all known
%Then we have the more suggestive. No GWAS is clear (lack of direct is okay for GWAS because direct is unbiased). For rare, lack of direct is trickier. Make point of disease relevant != disease associated (give example). Maybe we next go lack of direct for GWAS, then lack of direct for rare.
%Then this illustrates the two axes. On the other side is "novel". These are harder to interpret for GWAS. For rare, they suggest genes with interesting biology.
%Putting all of these together, we can think of a fingerprint. Where diagonal indicates well mapped biology, x-axis for 

Each relationship in the gene-trait map is a genetically estimated probability of genetic support combining both human genetic (direct) and non-human-genetic (indirect) data sources\ref{fig:direct_vs_indirect}. The 72,387 gene-trait pairs with probabilities >99\% are almost entirely well-established gene-trait relationships -- some of the strongest include \textit{GATA1} and rare hematologic disease, \textit{RPGR} and retinal ciliopathy, \textit{IL10} and inflammatory bowel disease, and {\it GCK} and blood glucose levels. All of these relationships are supported by both genetic associations (i.e. curated variant segregation evidence for \textit{GATA1} and \textit{RPGR}, strong GWAS signals nearby \textit{GCK} and \textit{IL10}) and extensive non-genetic evidence from the literature and biomedical databases (e.g. known biological process for \textit{GATA1}, mouse data for \textit{IL10}, and expression patterns for RPGR; \ref{supptab:examples_provenance}). 

Fig 1a: Direct vs. indirect across all gene-trait pairs, colored by trait category. Label examples 

%1.	Phosphate – Serum phosphate variation has known associations but little causal mechanism work via GWAS.
	%2.	Monocyte count – GWAS hits exist, but mechanistic links to disease pathways are thin.
	%3.	Eosinophil percentage of granulocytes – Mostly descriptive associations; mechanistic understanding is limited.
	%4.	Blood metabolite measurement (non-specific GCAT) – While some metabolites are well studied, many have uncharacterized causal pathways.
	%5.	Blood protein measurement (non-specific GCAT) – Similar to metabolites, many associations lack mechanistic follow-up.
	%6.	Gut microbiome measurement – GWAS here is very early-stage; causal biology is unclear.
	%7.	Lymphocyte percentage of leukocytes – Genetic influences known, but little mechanistic explanation in immune regulation.
	%8.	Body height (non-GCAT) – Extremely polygenic, but individual mechanistic pathways beyond growth factor biology are poorly resolved.
	%9.	Eosinophil count – Some asthma/allergy overlap is known, but causal networks are incomplete.
	%10.	Systemic lupus erythematosus (non-GCAT entry) – GWAS hits exist, but full mechanistic pathways for many hits remain undefined.

    %1.	Total body lean mass – Associations exist, but causal muscle biology from GWAS is thin.
	%2.	Bioavailable testosterone levels – Genetic links are known, but mechanistic endocrinology is poorly resolved.
	%3.	Waist-hip ratio adjusted for BMI – Polygenic patterns found, but underlying pathways beyond adipose biology remain vague.
	%4.	Cortical surface area measurement – Strong neuroimaging GWAS results, but no clear cellular mechanisms.
	%5.	Heel quantitative ultrasound estimated bone mineral density (eBMD) – Signals found, but bone microstructure biology remains unclear.
	%6.	Pulse pressure – Cardiovascular genetics has hits, but exact vascular mechanisms are elusive.
	%7.	Myocardial infarction – Known associations, but mechanistic mapping for many loci is lacking.
	%8.	Phosphate – Metabolism GWAS exists, but mineral regulation pathways from variants are underexplored.
	%9.	Primary biliary cholangitis – Immune-mediated liver disease, but for most loci, the mechanistic liver biology is unknown.
	%10.	eGFRcreat (serum creatinine) – Kidney function loci are known, but causal nephron mechanisms are incomplete.

Beyond these near-certain gene-trait relationships with PP>99\%, we find 163,317 ones with weaker but still strong genetic support (PP > 80\%). For common traits, many of these result when genes nearby GWAS hits are identified based on their strong indirect support relative to neighboring genes. For example, {\it EYA2} has strong support (PP=0.91) for Waist-hip ratio adjusted for BMI (an indicator of body fat distribution), due to its proximity to a strong GWAS association (rs4810576, p=$1.4 \times 10^{-50}$) and its indirect support from gene annotations\ref{supptab:examples_provenance} reflecting its known role in developmental processes\cite{tadjuidje_eyes_2012}, capturing a previous hypothesis that {\it EYA2} influences adipocyte differentiation or depot-specific gene regulatory programs\cite{perrin_chromatin_2021}. Two examples of relationships for Phosphate (a systemic marker of mineral metabolism) are \textit{FOXO2} (rs9486916; p=$4.4\times10^{-40}$) and \textit{AHR} (rs4410790; p=$1.4\times10^{-13}$), both of which have high indirect support due to  links to abnormal circulating hormone levels and abnormal primordial ovarian follicle morphology in mice\cite{essaghir_transcription_2009,hernandez-ochoa_ability_2010}, likely reflecting an endocrine-ovarian regulatory axis that may contribute to sex- and age-related differences in mineral metabolism. 

Many other complex trait relationships are for genes that do not lie near significant GWAS associations, such as \textit{GATA3} (PP=0.84)) and \textit{NFIB} (PP=0.53) for Chronic Obstructive Pulmonary Disease (COPD). Both lie near suggestive COPD associations (rs183304223, p=$3.7\times10^{-6}$ for GATA3 and rs877570, p=$9.6\times10^{-8}$ for NFIB) and have gene annotations for narrow eye opening annotations in mice\cite{kurek_transcriptome_2007} (in the case of \textit{GATA3}) or for lung development and morphogenesis\cite{abe_gata3_2021} (in the case of \textit{NFIB}), reflecting roles in epithelial structure and lung maturation pathways. There are also examples of relationships with a complete absence of direct genetic support: XXX have combined PP>0.5 but direct PP<X. Examples include \textit{CD19} for systemic lupus erythematosus (SLE), predicted because annotations for B-cell receptor signaling and autoantibody production capture its role in humoral autoimmunity\cite{mackensen_anti-cd19_2022}, and \textit{FGF8} for snoring (adjusted for BMI), predicted because annotations for craniofacial patterning and skeletal development capture its role in setting jaw/palate morphology\cite{zbasnik_fgf8_2022} and hence the size of upper-airway passages. [[[These examples are notable because GWAS are unbiased and well-powered, reflecting some error in our annotation of this gene or its properties]]

For rare diseases, we find XXX relationships without recorded direct genetic support in the Orphanet database. Relationships between \textit{RPL38} and orofacial clefting syndrome (PP=0.2, indirect support from links to cleft palate in mouse\cite{kondrashov_ribosome_2011}, co-occurence with Diamond-Blackfan Anemia in the literature, and interaction with EFTUD2 in the STRING database; \ref{supptab:examples_provenance}); \textit{CLRN2} and rare deafness (PP=0.2, indirect support from annotations regarding auditory brainstem malformations in mice\cite{mendia_clarin-2_2024} and involvement in the sensory perception biological process in the Gene Ontology;\ref{supptab:examples_provenance}); or {\it MESP1} for Tetralogy of Fallot (TOF) (PP=X, indirect support from annotations from cardiac septation and outflow-tract defects in mice\cite{lescroart_defining_2018}) are not considered to have strong clinical support but are potentially worth including as novel, mechanistically plausible candidate genes in diagnostic pipelines. As discussed in \ref{sec:methods}, the indirect support values for these genes model their disease-relevance, which is only one aspect (alongside mutational constraint and disease-pathway centrality) that determines whether they may harbor Mendelian disease-causing mutations.

Although both direct and indirect support are derived from human genetic data, they emphasize different data sources and can thus be thought of as two quantitative ``axes`` on which gene-trait relationships can be categorized. In Figure \ref{fig_direct_indirect}, genes close to the y-axis (no indirect support) are ``novel'' in the sense that the genes do not share annotations with known genes for the trait. Among the 1,096 common traits with power to detect indirect support (\ref{sec:methods}), the 19,260 gene-trait relationships with no indirect support but PP>50\% are likely highly enriched for GWAS associations for which experimental characterization could lead to new insights into disease -- for example, for type 2 diabetes (T2D) these genes include some with few annotations and about which little is known (\textit{MTMR8}, number of annotations in the 5th percentile across genes), together with some heavily annotated but which appear to act through biological processes distinct from those that charaterize the bulk of genes associated with T2D (\textit{GINS2}, number of annotations in the 90th percentile). For GWAS, such genes must be interpreted probabilistically, as uncertainties in linking SNPs to genes implies that the true effector gene may in fact have indirect support (for both \textit{MTMR8} and \textit{GINS2}, no genes within 1Mb of the lead GWAS SNP have indirect support, but this still does not rule out the possibility of even longer-range interactions\cite{claussnitzer_fto_2015}).

For rare diseases, genes with direct but no indirect support are ``novel'' (despite their appearance in curated gene catalogs) in the sense that they likely have functions different from those of other curated disease genes. Across the 210 rare diseases with power to detect indirect support, we detect 52 such genes across 22 diseases. For example, {\it PUS7} for non‑specific syndromic intellectual disability affects neurodevelopment through RNA modification and tRNA pseudouridylation, distinct from most other known intellectual‑disability genes that act through synaptic or transcriptional regulation\cite{shaheen_pus7_2019}; {\it DAAM2} for steroid‑resistant nephrotic syndrome affects podocyte function through actin cytoskeleton remodeling, distinct from most other nephrotic syndrome genes that act through slit‑diaphragm or ion‑channel complexes\cite{schneider_daam2_2020}, and {\it SCAPER} for Bardet–Biedl syndrome affects ciliogenesis through cell cycle and ciliary length regulation, distinct from most other BBS genes that act through BBSome or intraflagellar transport machinery\cite{wormser_scaper_2019}. Such genes  are of particular interest because they highlight alternative, underexplored disease mechanisms that could expand our understanding of rare disease pathogenesis.



Idea here: this map departs from the traditional significance threshold approach. Rather than a binary FDR above or below 0.1, we can read out entire "fingerprints". For example, when XXX is considered, the pattern of traits on the diagonal are ones in which our genetics matches our known biology, lending credence to its involvement in XXX even if some traits (i.e. XXX, XXX) have only moderate PPs. Additional indirect support values with no direct support can provide examples of mechanisms that are not supported genetically. For common traits, this is notable since GWAS are unbiased. For rare traits, the explanation could be that while the gene is disease relevant, it is not central enough to the trait to harbor monogenic mutations. Considering the full pattern of traits provides a statistical picture of the gene's function and relation to traits that can be queried for various questions (see agentic AI)

Finally, when considered across all traits for each gene, the pattern of direct and indirect support values provide a "fingerprint" that reflects how well our current expectations of gene function matches realized genetic associations for it. Genes in which indirect support correlates with direct support for the most part have established mechanisms. Example: MT-N5 for rare, APOC1 for common. [Pick a gene with a diagonal and then look at strong indirect with no direct for GWAS, or strong direct with no indirect for rare]. Pick two genes, one with peaked distribution and one with broad (PPARG?)
PDE3A
COL1A1
Many of the extreme off-diagnoal are associated with GWAS lipids -- need to run NMF?

Genes in which indirect support is uncorrelated with direct indicates genes whose known biology does not reflect their effects on human phenotypes: examples include COX18 and 

Some of this could be due to uneven sampling/power of traits (see Methods -- i.e. genes that affect metabolic diseases will have higher total sums, and maybe some trait areas have higher or lower power even if they have same number of GWAS)

1. Lies along the diagonal: biology of gene and relationship to disease is well understood (minimal number off diagonal)2. A lot of indirect but no direct: must be constraint no? What else (ChatGPT is answering) (not many on diagonal)3. A lot of direct but no indirect: maybe an error in GWAS mapping, or annotations gaps (not many on diagonal)4. Some traits on diagonal, but a fair number of others off (what does it mean? For those on x and those on y?) that it has a lot on diagonal means we understand it's role in disease, but then there are others that don't fit the pattern (maybe high indirect means pathway not fully relevant to disrase, high direct may indicate less well understood pleiotropy)5. Some traits high indirect and some high direct but none really on diagonal (what does it mean)? These are genes that whatever we think we know about the gene doesn't really capture its relationship to disease
So we need metrics to capture (a) fraction on the diagonal (average departure from 45) (b) fraction on y (c) fraction on x (d) direct minus indirect

Plot fraction on diagonal vs avg direct minus indirect

Can bucket x axis into three areas, left all concentrated on y axis, then spreads out going to right. Far right has three bins (direct only, indirect only, mixture of both). Middle has spread where gene is mostly known but there are some that are either acting through a different mechanism then thought (y) or don't actually contribute to disease (x)

Fingerprint of relevance from GWAS can predict rare disease -- see BSCL2


Fig 1b: plot of indirect sum vs. direct sum

Fig: phewas for a gene

Pick figure of gene with a wide spread of points across direct vs indirect plot


These relationships also highlight genes with interesting patterns across traits. Pattern across traits for a gene can be interesting. For example, XXX. For example, COL1A2 has indirect support for many traits, most of which have no direct support. This is in part explained by its high level of constraint, supporting it as a gene at the center of many biological pathways but intolerant to genetic variation. However, a handful of traits do show direct support, suggesting that mutations that impact human phenotypes are possible. Why these traits can be linked to perturbations, but others have not yet (or cannot be) could help understand the function of this gene. Conversely, PDE3B has direct support for many traits. but none with any significant indirect support. Genes such as this are clearly important physiologically, but are ``dark'' in the sense that we don't understand much about them. In some cases, the prediction may indicate that the gene is too constrained, for example ACTA1. In others, it is a bit of a mystery, such as PDX1. 

Fig: phewas direct vs. indirect for a gene Make beta vs beta marginal across all gene sets and all traits for a single gene

\subsection{Gene set-trait relationships that contribute to indirect support}

A secondary output of PIGEAN is the effect sizes that each gene annotation has on a gene's indirect support. When these effect sizes are estimated marginally (that is, not accounting for the effect sizes of other gene annotations with which it is correlated), the results can be viewed as sparse and numerically interpretable gene set enrichment results. For example, for CAD, the top gene sets are XXX. For sleep, they are XXX. For XXX, they are XXX. These values not only identify pathways that explain genetic predisposition to a trait.

Fig: joint vs. marginal across all gene sets and traits

Some of these are novel. Including XXX. And XXX. And XXX. Pick some from each library here. Therefore, even though gene sets capture known biology, they can reveal novel insights if that mechanism has not been linked to a trait previously.

Gene annotations are not synonymous with mechanisms, but in many cases they can be thought of as a noisy readout of a mechanism. For example, we can look at pathway XX and do phewas. Others such as mouse phenotypes are less mechanism, but still provide clues (e.g. insulin secretion vs. response). Mapping free text to gene sets is hard, but it is also possible to do mechanism first searches, particularly with the assistance of an AI agent. Examples of searches for a mechanism:  https://chatgpt.com/share/689e660e-2534-800d-a2c6-99003c3d8b6d. Another example: XXX.

The joint effect sizes allow us to find independent mechanisms. For example, with CAD, the top joint annotations are XX, XX, and XX, pointing to distinct mechanisms of lipids, vascular traits, and heart structure. The genes with these annotations and high direct support (XXX, XX, and XXX) are candidates to act through distinct mechanisms. We observe many different betas for rare diseases as well. For example, XXX has XXX. XXX has XXX. Rare diseases therefore also have evidence of acting through different pathways. Look across multiple gene sets with beta for a disease (mechanistic heterogeneity)
 - search which traits have heterogeneity clinically but not clear molecularly
 - do for both rare and common

In many cases the joint and marginal values are quite different. There are far more marginally significant betas than there are jointly significant. For XXX. For XXX. In many cases, gene sets that might appear to be predictive are not after accounting for other gene sets. For example XXX, which is captured by XXX. Or XXX, which is captured by XXX.

Established gene annotations does not mean only known genes are found, since genetics picks out which ones matter. Note the indirect support estimated from the narrow eye opening mouse annotation suggests how genetic calibration can identify potentially unexpected disease relationships even from established gene annotations.

\subsection{Mechanistic insights into gene-trait relationships}

In addition, gene sets also help assign a mechanistic explanation for a gene-trait explanation. For example, XXX has an association for X, and gene sets suggest XXX. And XXX. These can be interpreted with the help of an AI agent; passing it the gene, trait, and the gene sets can suggest a hypothesis.

In most cases, most annotations of a gene are irrelevant to the trait. For example XXX has XXX annotations, only X of which have non-zero marginal effects and XX have joint. The vast majority are irrelevant, which significantly simplifies the challenge of understanding the gene function. But others are surprisingly irrelevant, such as XXX. Other examples include XXX.

Novel genes for a trait of interest can also be interpreted based on their indirect support for other traits. For example, novel genes for XXX include XXX (all have no indirect support). However, each gene has indirect support for other traits. XXX has XXX, suggesting XXX. XXX has XXX, suggesting XXX. Therefore, the numerical vector of indirect support values across traits provides mechanistic information about a gene that can be used to uncover new mechanisms for a primary trait of interest.

There are also examples, however, of genes that appear to affect different traits through different mechanisms. For example, PPARG to T2D is different from PPARG to bone disease. Other examples include XXX. These are genes that could belong to different pathways, as captured by distinct gene annotations, and the different pathways are linked to different traits. This might be relevant in understanding or avoiding drug side effects.

\subsection{Validation of the resource}

%1. Hold out cross validation to show indirect predicts direct
%2. Temporal to show combined predicts future
%3. Stratify across trait architectures and sample size, rare vs. common
%4. Bias towards well studied genes or traits with higher power
%5. Comparison to direct only baseline

% Figure 2 panels
%a. Show indirect predicts direct in hold one out chromosome (portal_model)
%b. Show mechanistic alignment in portal model results compared to null for combined
%c. Show relative success results compared to benchmarks for combined (portal model)
%d. Show it predicts future GWAS signal (portal) 


\begin{figure}[t!]
    \centering
    \includegraphics[width=1\linewidth]{figures//indirect-support/validation-figure-1}
    \caption{(a) A subset of pigean models run on a 13 trait common disease validation set. Displaying top performing models and models selected through assessment of top genesets. Marker size represents the complexity of the model and the multiple transparent grey boxes show overlayed null distrbutions in terms of both relative success and mechanistic alignment for each model (Methods). (b) Same models run over a 40 trait rare disease valiation set. (c) Over all model selection experiments the proportion of genesets that came from each geneset library that where significant for a given trait-model pair. }
    \label{fig:model-selection-main}
\end{figure}

% General overview paragraph of validation
We next sought to systematically validate PIGEAN along two goals: (i) predicting genes that are biologically and clinically meaningful and (ii) identifying gene sets that are mechanistically aligned with each phenotype. To do so, we evaluated a focused benchmark comprising $n=13$ common traits selected for representative clinical coverage and statistical power, and $n=40$ rare traits partitioned into $n=20$ disease-specific lists and $n=20$ disorder-category lists (see Methods). Because our full universe spans $>7,000$ traits, an exhaustive search over the model space was computationally infeasible; instead, we executed a targeted model-selection procedure that iteratively pruned candidates while answering specific research questions about PIGEAN’s behavior. For model selection, we defined a Minimum Description Length (MDL) criterion \cite{} that balances complexity and accuracy, where complexity captures both the number of gene sets made available to the model and the number of gene sets deemed significant for prediction, and accuracy integrates (a) potential clinical success of nominated targets and (b) mechanistic alignment between significant gene-set semantics and the phenotype under study. Operationally, this yielded three independent axes for comparison across models: (1) clinical success, quantified as the relative success \cite{minikel_refining_2024} of PIGEAN-prioritized genes that exceed a threshold of direct/indirect genetic support; (2) mechanistic alignment, quantified as the normalized discounted cumulative gain (nDCG) of semantic similarity ranks between significant gene-set names and their corresponding phenotype; and (3) parsimony, quantified as the balance between the number of predictive gene sets surpassing a $\beta$ threshold and the total size of the gene-set libraries included.

% Detailed description of top performing model with pointers to supplementary information. I think it helps ground the reader in the analysis when discussing just a single models performance
The top-performing model across the common and rare validation sets used two libraries: MSigDB \cite{} (with HPO-derived gene sets removed to avoid circularity) and 2024 Mouse Phenotypes \cite{}, for a total of $41,627$ gene sets. Under this configuration, the relative success for common disease was 8.623 (95\% CI: 5.499–13.534), representing a near 4-fold increase over prior reports \cite{minikel_refining_2024}, and for indirect-only genetic support it was 3.255 (95\% CI: 1.748–6.059), a regime that, to our knowledge, has not been previously characterized. Further, from rare traits using only a curated set of rare gene lists, we achieved a relative success of 2.167 (95\% CI: 1.028-4.57) which is only from indirect-support only as direct support in this context leverages the rare disease gene lists which would inherently be circular. Owing to PIGEAN’s regularization, these clinical-success gains were driven by an average of $124$ significant gene sets per common trait and $139$ per rare trait (Supplementary Information), amounting to approximately $0.3\%$ of the $41,627$ available gene sets. It should be noted that these relative success results are only amoung the validation traits in this analysis, and upon analyzing the full phenome of $>7000$ traits which is reported in the subsquent sections.

Mechanistic alignment, measured via nDCG (idealized ranking $=1$), was $0.79\pm 0.09$ , and the average phenotypic similarity of the top $k=10$ $\beta$-ranked gene sets was $0.202$ for both common and rare phenotypes. Finally, increase of phenotypic similarity among significant gene sets over a random null was statistically significant for common traits ($p=6.36\times 10^{-19}$) and for rare traits ($p=1.36\times 10^{-5}$) demostrating PIGEAN is capturing a genetic signal that is supported by literature. Figure~\ref{fig:model-selection-main} shows a subset of evaluated models that are of particular interest, and it is visually clear that the MsigDB + Mouse Phenotype model performs well overall balancing clinical relevance, mechanistic alignment, and model complexity. 

\begin{figure}[t!]
    \centering
    \includegraphics[width=1\linewidth]{figures//indirect-support/validation-figure-2.png}
    \caption{Relative success of drug targets launched post-2018 over baseline Magma supported genes using non-cicurlar models using gene sets curated from pre-2018 sources priorized (a) with direct genetic support and (b) indirect genetic supprt with secondary $y$-axis showing number of predicted targets for each genetic support threshold. Improvement in regression fit ($\Delta R^2$) of regressing Magma Z-scores from GWAS pre-2018 on GWAS post-2018 with (c) a direct genetic support filter and (d) an indirect genetic support filter (outlier traits removed).}
    \label{fig:validation-figure-2}
\end{figure}

% Back to the future validation
Our final validation step was designed to assess the potential impact of circularity between GWAS results and curated geneset libraries in predicting drug targets as well as in forecasting future GWAS significance. Because many published genesets are curated from GWAS signals, we sought to confirm that PIGEAN was not simply reinforcing known signals but was instead capable of highlighting novel associations. Similarly, we repeated our drug target relative success analysis using only targets that would not have been available prior to the publication of a given geneset. To this end, we created 10 additional models using gene set sources from MsigDB, Mouse Phenotypes, and MeSH Terms. For each library, we curated pre-2018 and post-2018 versions based on provenance information. We then applied our validation analyses across 67 traits with both pre- and post-2018 GWAS. For the relative success analyses, only drug targets launched after 2018 were included to avoid target circularity.  

As shown in Figure~\ref{fig:validation-figure-2}a-b, we observe consistent improvements in relative success as a function of both direct and indirect support, even after temporally excluding any form of GWAS or drug target information. Notably, the pre-2018 text-mining model, which explicitly removed all GWAS and drug MeSH terms (Methods), achieved the strongest performance. At a direct genetic support log-odds threshold of 4 ($P=98.2\%$ relevance), this model reached a relative success of 4.51 (95\% CI: 2.67–7.63) across 50 predicted targets, significantly above the Magma baseline ($p$-value $=2.5\times 10^{-6}$). Although somewhat weaker, indirect genetic support also improved relative success, reaching 3.04 (95\% CI: 1.93–4.80) at the same log-odds threshold with 108 predicted targets. This highlights an important distinction between Figure~\ref{fig:validation-figure-2}a and~\ref{fig:validation-figure-2}b: indirect support identifies additional targets not captured by direct genetic support alone, as reflected in the plotted trend lines showing the number of predicted targets as a function of the (in)direct support thresholds. We further analyze the complementarity of direct and indirect support in expanding drug target diversity in the following section.  

We next evaluated how well direct and indirect genetic support could predict future GWAS significance by regressing on post-2018 MAGMA Z-scores. When directly regressing Z-scores (Methods), we observed modest but statistically significant gains in model fit ($\Delta R^2$) when incorporating indirect support, with maximum improvements across all traits of rangeing from 0.0019 to 0.017 (Extended Data Figure XXXX). The proportion of traits with significant $\Delta R^2$ also varied by model (ranging from 62.9\% to 91.4\%). However, substantial noise remained in non-significant genes, where MAGMA Z-scores were nearly random, such that even baseline regressions of pre-2018 versus post-2018 Z-scores yielded a low average $R^2 = 0.14$ across 67 traits.  

To test whether genetic support could act as a filter to reduce noise and forcast genes that will be predictive of future GWAS, we regressed post-2018 MAGMA Z-scores on pre-2018 scores while progressively restricting to genes meeting increasing thresholds of direct or indirect genetic support (Figure~\ref{fig:validation-figure-2}c-d). Across all models, regression slopes remained positive for both direct ($\beta$ range: 0.014–0.038) and indirect ($\beta$ range: 0.050–0.058) support, with indirect support providing a greater increase in model fit. Importantly, filtering on stronger evidence not only increased average $\Delta R^2$ but also strengthened the concordance of pre- and post-2018 associations, consistent with the notion that PIGEAN prioritizes genes with reproducible relevance. Together, these temporal validations support that PIGEAN is not simply recapitulating latent GWAS or drug target information embedded in curated libraries, but instead enriches for signals that are predictive of future discoveries.  

\subsection{Global properties of gene-trait relationships}

"Our examples show that the global map can recapitulate manual searching. Its real power though comes from it being complete and quantitative. This motivates these landscape analyses..."

Plot of sum of betas across traits
Plot of number of genes with high total vs. number with high indirect (points are traits); also number of genes with high indirect vs number with high direct
How many greater than 50\%? How many unique genes of these?
How many greater than 50\% on average for a trait? Use 5\% threshold for common and 0.5\% threshold for rare

Make plot of sum direct vs sum indirect.
 - the sum of indirect is in general larger
 - common traits have correlation more strongly with pLI than rare (is this a power issue though)
 - pLI correlates with sum of indirect and sum of direct, but direct more strongly (confounders in trait ascertainment, confounders in power may be an issue here)
  - for genes like APOC4, a lot of traits with no indirect. But how much of that is due to lack of power of the GWAS? We can't say necessarily that those traits are novel findings for it


Amount of redundant information in gene set libraries -- plot sum of betas vs. sum of betas uncorrected. The difference indicates how many gene sets are actually redundant.

Sum of betas for rare vs. sum of betas for common -- any more or less coherence 

Plot number of clumps vs sum of betas

Plot sum indirect or sum direct as a function of number of annotations (have to filter and correct for traits first)

How many effector gene predictions do we make
 - how many cases does probability without gene sets increase after with gene sets
 - how many genes near a lead SNP have probability above X

 Histogram of number of effective genes per loci (need the ge.out)

Investigate what lack of direct means for rare. pLoF/LOEUF vs. indirect-direct etc. The ones with indirect + direct likely more constrained, but indirect in general more constrained than non-indirect

Then what lack of indirect means for GWAS. Maybe here we have histogram of number of genes per locus with prior above X. Show what happens to direct with gene sets vs. direct w/o gene sets

\subsection{Calibrating non-human-genetic data sources across different human traits}

Which libraries are good/bad
Which libraries vary in utility across trait groups


\subsection{Prioritizing drug targets with indirect genetic support}

In addition to RS, maybe a plot with (a) a histogram of the PIGEAN combined scores for approved drug targets (b) a second histogram of failed (c) the overall histogram across all genes (d) a cutoff for score based on the approved/failed to indicate how many gene-trait pairs we have in the resource that are likely to suceed as drug targets. Then statistics on how many of these are for diseases with no treatments, how many of these are for genes with drugs for a different disease (repurposing), how many of these are for genes in different IDG studies (dark/clin/bio etc.). Then give a few examples of potential targets in each category

\subsection{Analyzing gene-trait relationships with agentic AI}

"Because this map is dense we can use it in ML..." then do that

Guess rare diseases from common trait vector (combined, direct, indirect). Give examples.

\section{Discussion}

State limitations that is statistical and integrating available data sources rather than new data. Need to very clearly state it is a map, a resource. Our analyses show how much it can help in hypothesis generation. Confirmation requires experiments. But, it gives a starting point by filtering to high confidence genes, and giving mechanistic clues based on the pattern of associations across all traits as well as the gene annotations that explain the indirect genetic support.

Globally, the statistical properties are clear. RS, correlations, calibration, etc. So distinguish examples from statistics.

Discuss how this builds up on PoPS with key advances: interpretable, statistical, accurate, combine with direct, applies to rare

Also in discussion compare to LLMs to prioritize genes using non genetic (statistical, quantitative, searchable, for ML) -- this may be a bigger issue for grouping with EAGGL

And KGs (dense, embedding, statistical)

Where does this fit in regarding gene set completion (GSFM)




\section{Older text}
(Fig.~\ref{fig:pigean-concept})
% Introduce PIGEAN and how it instantiates the concept of indiret genetic support efficiently
We define indirect genetic support as the probabilistic influence that genetically implicated genes exert on related genes through shared (functional) annotations. These annotations, such as gene sets from , provide structured, non-genetic information that links genes based on common biological contexts (Fig.~\ref{fig:pigean-concept}). While indirect support could in principle be modeled in various ways, our approach formalizes it as an informed prior (\hyperlink{pigeanpior}{Methods}) that intends to capture more prior relevance if a gene is in one or multiple shared contexts that have other genes with high direct genetic support for a trait. For each gene set library, we construct a matrix, $X$, that encodes gene–annotation relationships, and estimate the contribution of each annotation to trait relevance, $\beta$, using a sparsity-enforced logistic regression model (\hyperlink{pigeanpior}{Methods}). We enforce sparsity by sampling the hyper-prior distribution over $\beta$ as a spike-and-slab distribution \cite{mitchell_bayesian_1988,madigan_model_1994,ishwaran_spike_2005} (\hyperlink{pigeanpior}{Methods}, Extended Data Fig~\ref{}). We find that biologically related traits tend to prioritize similar gene sets, supporting the model’s ability to recover coherent and interpretable biological structure (Fig.~\ref{fig:geneset-cluster}).

% Explain direct support calculation
We employ multiple methods for modeling direct genetic support based on the type of genetic data available for a phenotype. If GWAS is available, we use an extension of a previously proposed Bayesian model \cite{dornbos_evaluating_2022} that retains simplicity and speed but includes extensions to account for suggestive association signals, the number of genes nearby a SNP association, association fine mapping, and probabilistic SNP to gene linkage functions (\hyperlink{pigeanpior}{Methods}). For rare disease associations, we use manually annotated probabilities of association derived from public databases. Additional extensions to complex trait exome sequence associations or rare disease de novo or case control Bayesian models \cite{TADA} are straightforward (\hyperlink{pigeanpior}{Methods}) but full analysis of these generatlizations are left to future work.

% Explain the model parameters and fitting using the iterative gibbs sampling
For each trait, we jointly fit this model across all genes and all gene sets through Gibbs sampling. In essence, this approach alternates between updating the gene genetic support estimates and the annotation effect estiamtes (\hyperlink{pigeanpior}{Methods}); in this way, genetic support estimates are refined as information passes through annotations shared by genes. It is this perspective that motivates the definition of ``indirect genetic support" (the contribution to genetic support from gene annotations and their effects) as an extension to the traditional notion of direct genetic support (the contribution from the genetic data for a gene).  At convergence the Gibbs sampler yields (a) jointly estimated effects on indirect support across all annotations in the model, and (b) combined genetic support values for all genes in the genome. To dissect these results, we also record (c) the marginal effects on indirect support from each annotation (i.e. not controlling for other annotations in the model, a value analogous to the traditional notion of gene set enrichment \cite{?}) and (d) a decomposition of the combined genetic support for each gene into direct and indirect values. The full model is depicted as a probabilistic graphical model in Fig.~\ref{fig:pigean-concept}C. 

% Explain the model parameters and fitting using the iterative gibbs sampling
% To jointly estimate both direct and indirect genetic support, we fit a model that alternates between updating the relevance of each gene to the trait and refining the functional annotations that best explain those relevance patterns (\hyperlink{pigeanpior}{Methods}). Specifically, given an initial estimate of gene relevance based on direct genetic support, we infer which gene sets are most informative by learning their contribution to explaining these relevance scores (logistic regression). We then use these refined gene set weights to update our estimates of gene relevance, now informed by both genetic and annotation contexts. By iterating these two steps by refining annotation weights and updating gene relevance, we converge on a joint distribution that captures both direct genetic signals and their propagation through structured biological knowledge. This inference procedure is implemented efficiently via Gibbs sampling, and the full model is depicted as a probabilistic graphical model in Fig.~\ref{fig:pigean-concept}.

\newpage

% Old Sections
\section{Results}\label{sec:results}


\subsection{Addressing Circularity Concerns Between Gene Annotations and Genetic Data}

%JF COMMENT: This is great and a significant contribution. But is it true? When in the past I ran PIGEAN with the gene sets mined from the literature that included genes from T2D GWAS abstracts, the model converged to have those gene sets as super highly weighted and biased significantly towards those genes at the top. So the model isn't squaring with the observation I had...do we need to update PIGEAN's inference approach somehow to account for this?

A common concern in the genetics community is the potential \emph{circularity} introduced when incorporating curated gene annotations (e.g., pathways, enrichments) alongside genetic data (e.g., GWAS, WES) in statistical models. This concern stems from the fact that both sources may reflect overlapping biological signals. Moreover, many gene annotations may themselves have been constructed with knowledge of prior genetic findings, such as GWAS hits, and potentially resulting in a double-counting effect when used together in gene prioritization models and confounding downstream analyses.

However, as illustrated by our proposed causal diagram (Figure~\ref{fig:pigean-causal-model}a), this circularity is not intrinsic if the underlying causal relationships are correctly specified. In particular, if both annotations ($X$) and genetic data ($S$) are treated as noisy observations of an unobserved biological state ($U$), and we jointly condition on both in a properly structured model, then the use of both sources does not induce circularity, even if gene annotations were partially informed by genetic data. 

We demonstrate this in a toy example (Figure~\ref{fig:pigean-causal-model}b) using simulated structural equations (see Methods) where the task is to recover the ground truth relevance vector $d \in \{0,1\}$ across an arbitrary set of genes $g_1, \dots, g_N$. Even as the correlation between $S$ and $X$ increases, simulating stronger causal influence of genetics on annotations, our model’s predictive performance (AUC) remains stable. The only observed effect is an inflation in the standard error of the effect size associated with $S$, which impacts classical p-value calculation but not the Bayesian posterior. In a Bayesian framework, this uncertainty is naturally integrated into inference, mitigating its impact.

%JF COMMENT: I think we need another paragraph to set the framework for how we are getting from the framing in the introduction (gene-trait relationships that incorporate two types of data) to this prior framing. The paragraph should start with the goal, state the high level approach (BF + prior is the two ways), pose questions as to how to learn the prior, how annotations fit, etc. In other words, we should have a paragraph that explains the intuition without having to read any math

Our method, PIGEAN (\textbf{P}riors \textbf{I}nferred from \textbf{GE}ne \textbf{AN}notations) implements a joint Bayesian model aligned with the causal structure in Figure~\ref{fig:pigean-causal-model}a. For a trait $T$ and genes $g_1, \dots, g_N$, let:
\begin{itemize}
\item $d_i \in \{0,1\}$ denote whether gene $g_i$ is causally relevant to the trait,
\item $\mathbf{X}_i = \{X_{ij}\}$ denote the set of annotations for gene $g_i$ (e.g., gene set memberships, pathway inclusion, expression scores),
\item $S_i$ denote the genetic evidence for $g_i$ (e.g., effect size, summary p-value).
\end{itemize}

The key innovation in PIGEAN is its use of $\mathbf{X}_i$ and learned weights $\beta$ to define an informative prior over $d_i$ via a logistic link:

\begin{equation} \label{eq:priors}
\log\frac{y_i}{1-y_i} = \sum_j X_{ij}\beta_j + \alpha_i + \log\frac{\pi}{1-\pi}
\end{equation}
\begin{align*}
d_i \sim \text{Bern}(y_i) \quad
\beta_j \sim \delta_j\mathbf{N}(0, \sigma_j^2) \quad
\alpha_i = \alpha\sum_j X_{ij} \quad
\sigma_j^2 = \sigma^2\text{Var}_i(X_{ij})^w \quad
\delta_j \sim \text{Bern}(p) \quad
\end{align*}

This construction allows us to model $d_i$ using a ``principled prior" derived from functional annotations, mirroring how domain experts reason about gene relevance using prior biological knowledge. While this prior depends on $X_i$ and $\beta$, it does not violate our causal model: in Bayesian inference, a misspecified prior can be corrected by the likelihood, and uncertainty in the prior is naturally incorporated into the posterior.

We then incorporate the genetic score $S_i$ as independent, complementary evidence for $d_i$. Since $S_i$ is a child of $d_i$ in the causal graph, we can model $P(S_i \mid d_i)$ (see Methods). The posterior distribution over $d$ and $\beta$ factors naturally via conditional independence:

\begin{equation}
\underbrace{P(d, \beta \mid S, X)}_{\text{joint posterior}} \propto 
\underbrace{P(S \mid d)}_{\text{direct evidence (likelihood)}} \cdot 
\underbrace{P(d \mid X, \beta)}_{\text{indirect evidence (prior)}} \cdot 
\underbrace{P(\beta)}_{\text{hyperprior}}
\end{equation}

Note that $S_i$ is conditionally independent of $(\mathbf{X}_i, \beta)$ given $d_i$, and $X_i$ is conditionally independent of $d_i$ given $U_i$ (which is latent). While marginalizing over $U_i$ induces a dependency between $X_i$ and $d_i$, we model this relationship explicitly via the logistic prior in Equation~\ref{eq:priors}. Moreover, $\beta$ does not enter the likelihood $P(S_i \mid d_i)$ and plays no role in modeling the direct genetic evidence.

A full visual overview of the PIGEAN workflow is shown in Figure~\ref{fig:pigean-concept}. A complete model specification and inference details are provided in the Methods section.

\subsection{Validation on Simulations} \label{sec:validation-sim}
We validated PIGEAN’s behavior through a suite of simulations designed to assess both specificity and robustness under varying conditions. As expected, when applied to null GWAS datasets, PIGEAN correctly inferred annotation effect sizes near zero and returned negligible direct evidence, demonstrating no spurious signal inflation. This was also true when annotations were randomly permuted across genes, supporting the specificity of the model’s inference procedure. These results suggest that combined gene relevance incorporating both direct and indirect evidence and measured as an approximate Bayes Factor (ABF) greater than 3 (measured in log odds) and annotation effect sizes above 0.005 can serve as conservative thresholds for study-wide significance (Extended Data Figure~\ref{fig:null_sim}, Supplementary Table 3).

Further simulations highlighted the distinct effects of annotation noise and genetic association noise on model performance. Focusing on annotation noise, we observed four key properties: (1) PIGEAN conservatively estimates annotation effects and gene ABFs due to spike-and-slab shrinkage; (2) annotations with negative effects or weak evidence are more aggressively shrunk, particularly in settings with few disease-relevant genes; (3) accuracy improves with increased annotation effect size and sparsity level; and (4) genes with many annotations must be regularized to avoid ABF inflation from asymmetric standard errors. Importantly, model performance was robust to hyperparameter mis-specification, with even large deviations in $\sigma$ and $p$ (specified in Equation~\ref{eq:priors}) yielding stable inference accuracy. Simulations incorporating GWAS summary statistics further confirmed that while association noise can reduce power, it does not induce false positives. PIGEAN’s conservative estimates under both annotation and association noise affirm its reliability for gene prioritization tasks. For more information, see Methods.

\subsection{Prioritizing Genes Across 13 Complex Traits} \label{sec:validation-real}
While many existing methods prioritize genes using either direct genetic evidence or indirect functional information, few frameworks explicitly integrate both. MAGMA \cite{leeuw_magma_2015} ranks genes based on proximity to genome-wide significant variants, capturing direct association evidence. In contrast, PoPS \cite{weeks_leveraging_2023} uses polygenic enrichments of gene-level features (e.g., gene sets, expression) to prioritize genes based on indirect functional relevance.

%JF COMMENT: Suggested framing if results stay the same:
%Show pigean can reduce to pops, verify with correlation on depict (maybe just put this in methods)
%Then run on all gene sets and show is better, maybe adding additional metrics if needed
%Then show it works more similarly on fewer gene sets (depict)
%Lead that into the number of parameters section

To benchmark the performance of PIGEAN against these widely used approaches, we evaluated genome-wide gene prioritization across 13 complex traits. For each trait, PIGEAN and PoPS (as PoPS internally leverages MAGMA) were used to generate ranked gene lists, which we then assessed for concordance with established gold-standard genes lists for these traits [Cite LLM Paper].



We quantitatively compared prioritization performance using a non-parametric Mann–Whitney $U$ test across the 13 trait contexts, following Demšar's methodology for paired method comparisons \cite{demsar_statistical_2006}. Using DEPICT gene sets as a first reference, PIGEAN performed equivalently to PoPS under the null hypothesis of no difference ($U = 113$, $p = 0.297$). This equivalence held ($U = 86$, $p = 0.842$) when using a broader set gene sets (see Methods) curated from PoPS, Mouse Phentotypes, Orphanet, Human OCR data, MSigDB, MeSH topics, and String DB. Figure~\ref{fig:pops-vs-pigean-goldstandard-depict} and Figure~\ref{fig:pops-vs-pigean-goldstandard-all} show the precision and recall @$k$ performance as function of $k$. Note: These are old, results and may improve once validation is redone. 

These results indicate that PIGEAN achieves gene prioritization performance on par with state-of-the-art methods while offering a unified, interpretable framework that incorporates both direct and indirect signals in a principled Bayesian model.

\subsection{Enhanced Performance With Fewer Parameters}

%JF COMMENT: Is this a placeholder for assuming the results will change to PIGEAN clearly better? The figures look already like it is better on the full model even if not significantly better. I wonder also if we need to expand our positive control gene list using the new o3 model...I'll bet one issue with the numbers is the incompleteness of the gene lists used as a gold standard

To better understand where PIGEAN's performance gains originate, we decomposed the model's contributions into indirect and direct components, and evaluated model sparsity and parameter efficiency.

Compared to PoPS, PIGEAN achieved equivalent or better prioritization performance using significantly fewer features (mean number of non-zero annotation weights: XXXXfor PIGEAN vs. XXXX for PoPS). This reflects the benefit of PIGEAN’s spike-and-slab prior and joint inference mechanism, which encourages sparsity while maintaining statistical power. The iterative re-estimation of gene and annotation weights was critical—ablating this step reduced precision@100 from XXXX to XXXX (Supplementary Figure~X).

In terms of direct support, MAGMA captured genome-wide significant loci but produced diffuse rankings across the gene body. PIGEAN, by contrast, concentrated high scores on functionally coherent genes. Among the top 100 PIGEAN-prioritized genes per trait, a median of XXXX overlapped with rare variant hits, compared to XXXX for MAGMA. Further, annotations with joint effects identified by PIGEAN were sparser (XXXX) non-zero annotations per trait vs. XXXX in PoPS), yet achieved stronger enrichment for known gene sets (mean enrichment p-value: XXXX for PIGEAN vs. XXXX for PoPS; see Supplementary Table~X).

Together, these findings indicate that PIGEAN can produce interpretable, biologically plausible gene rankings with fewer parameters, provided the full joint inference procedure is employed.

\subsection{Predicting the Future from Past GWAS}
We next evaluated whether gene scores inferred from older GWAS could anticipate hits from future, better-powered studies as a proxy for robustness and generalizability. For each trait, we identified a pair of historical and recent GWAS, and ran PIGEAN and PoPS on the older dataset. We then calculated the overlap between the top-ranked genes from the older study and those deemed most significant in the newer GWAS.

PIGEAN's predictions were the most stable across GWAS vintages, recovering a median of XXXX\% of the top 100 genes from the newer GWAS across traits, compared to XXXX\% for PoPS. In particular, PIGEAN prioritized XXXX genes that were undetected in the older GWAS but later confirmed in the updated study (see Supplementary Table~X). This suggests that PIGEAN's use of prior biological knowledge improves power to detect relevant genes even when sample size is limited.

We also observed that PIGEAN produced more consistent gene rankings across GWAS vintages (median Spearman $\rho$ of XXXX across 13 traits), suggesting that its integration of stable prior information buffers against fluctuations in statistical power.

\subsection{Validation on 40 Rare Diseases} \label{sec:rare-disease-validation}

\begin{figure}
    \centering
    \includegraphics[width=0.5\linewidth]{figures/indirect-support/real/rare-disease-validation.png}
    \caption{Example rare disease validation curve. Precision@$k$ is computed across simulations where each disease-associated gene is held out and scored using indirect evidence from the remaining curated gene list.}
    \label{fig:rare-disease-validation-sample}
\end{figure}

To further evaluate PIGEAN’s ability to capture biologically meaningful gene–trait relationships using indirect evidence alone, we conducted validation across 40 rare diseases sourced from Orphanet. These included both well-defined disorders with curated gene lists from OMIM as well as broader phenotypic categories spanning multiple tissues and systems.

For each trait, we simulated a scenario in which all but one known associated gene had been observed. Specifically, we initialized each known causal gene with a prior probability of 0.95, and set all other genes to a background prior of 0.05. We then computed indirect PIGEAN scores using only the known genes (excluding the held-out one) and assessed the rank of the held-out gene in the resulting prioritization. This design mimics a use case where the user seeks to discover one final relevant gene based on indirect support from existing known biology.

As shown in Figure~\ref{fig:rare-disease-validation-sample}, PIGEAN rapidly accumulates predictive precision—achieving over 50\% hit rate within the top 50 genes on average, and recovering nearly all true positives within the top 500. These results demonstrate that PIGEAN’s use of gene annotations provides a powerful mechanism for prioritizing candidate genes even in low-sample or incomplete trait settings.

\subsection{Model Selection and an Assessment of Gene Annotation Quality} \label{sec:model-selection}

After validating PIGEAN’s gene prioritization performance on both common and rare diseases using curated gene lists, we next sought to understand how the choice of gene annotation libraries impacts performance. To address this, we first estimated the prior inclusion probability $p$ for each gene set library using a large collection of GWAS datasets, restricting the model to a single library at a time (see Methods). This allowed us to derive a global informativeness score for each library, capturing how often its annotations were inferred to be relevant across a variety of traits of different powers.

We then conducted model selection experiments by iteratively constructing candidate models of varying sizes, each containing different combinations of gene set libraries, including both high-confidence sources and more circular or noisy ones (e.g., text-mined sets that incorporate GWAS-derived terms). This enabled us to test whether PIGEAN's performance improves with curated, biologically grounded annotations or whether it can also leverage noisier sources effectively.

An example of this analysis is shown in Figure~\ref{fig:rare-full-model-genesets}, which highlights the number of gene sets from each library that appeared among the top 1000 annotations contributing to rare disease prioritization. These early results suggest that annotation quality and redundancy both play critical roles in model effectiveness, with some libraries consistently contributing more informative features. Note: will extend this analysis to common traits and full model selection procedures.

\subsection{Exploring a Public Resource of \getvar{var__NUM_PIGEAN_TRAITS} Common and Rare Phenotypes}
\textbf{Goal:} Present exploratory findings and infrastructure built to support hypothesis generation from PIGEAN outputs.

Outline:
\begin{itemize}
    \item Overview of Resource
    \begin{itemize}
        \item Number of traits (\getvar{var__NUM_PIGEAN_TRAITS}), split across common vs. rare
        \item Description of portal interface or access model
    \end{itemize}
    \item Key Findings / Case Studies
    \begin{itemize}
        \item Showing the direct vs. indirect table for traits
        \item Notable annotation or gene cluster (brief summary)
        \item Surprising indirect-only hits (brief summary)
    \end{itemize}
    \item The PGEG AI Agent Workflow
    \begin{itemize}
        \item Description of the agent architecture (input: trait + scores; output: narrative/mechanistic insight)
        \item Examples of successful insights or hypotheses generated
        \item Linking to other RAG systems or real-world evidence 
    \end{itemize}
\end{itemize}

Other ideas:
\begin{itemize}
\item How many gene/phenotype links are for genes not in a GWAS locus?
\item What fraction are indirect only?
\item What is the correlation? What fraction are direct only?
\item What trait pairs have high direct overlap but less indirect overlap? Vice versa?
\item How many GWAS loci are "solved" by gene sets? How many have more than one?
\item 
\end{itemize}




\subsection{Combining Direct and Indirect Support Increases Relative Clinical Success of Drug Targets}
By leveraging a curated public dataset of filtered Citeline Pharmaprojects monotherapy drug programs since 2000 \cite{minikel_refining_2024}, we were able to study PIGEAN-based gene priortization as a function of clinical success. The curated Pharmaprojects dataset contained \getvar{var__NUM_UNIQUE_DRUG_PROGRAMS} unique drug programs indexed by target-indication pairs where targets are genes and indications are traits defined by the Medical Subject Headings (MeSH) ontology. We mapped \getvar{var__NUM_PIGEAN_TRAITS} traits avaiable in our Association to Function (A2F) web resource to associated MeSH headings yielding \getvar{var__NUM_UNIQUE_MAPPED_TRAITS} mapped traits over which we could analyze PIGEAN target priortization in terms of clinical success. This lead to a total of \getvar{var__NUM_PIGEAN_DRUG_PROGRAMS} unqiue drug programs that had both minimal combined PIGEAN support ($ABF_c \geq 1$ or $\sim 73\%$ probability that a target is relevant to a trait), and a trait indication with at least 80\% semantic similarity with a respective drug program (See Methods). This led to a \getvar{var__ASSOCIATION_PERCENT_INCREASE_FROM_MINIKEL}x increase of avaiable drug programs with genetic support from previously curated genetic assocations from various sources \cite{minikel_refining_2024}. We chose our combined support threshold of 1, as this yeilded \getvar{var__NUM_DRUG_PROGRAMS_WITH_DIRECT_SUPPORT} programs with direct PIGEAN support which was comparable to the \getvar{var__NUM_DRUG_PROGRAMS_WITH_GENETIC_SUPPORT_MINIKEL} programs previously analyzed \cite{minikel_refining_2024} with a \getvar{var__MINIKEL_PIGEAN_DIRECT_SUPPORT_OVERLAP}\% overlap. 

In all analyses, relative success (RS) is defined as the ratio of the probability of success of drug programs moving to a subsequent phase (e.g., Phase I to Phase II) of programs that have genetic support versus programs that do not have genetic support. In the case of PIGEAN, this translates to having a target with direct, indirect, or combined support (depending on the analysis) beyond a specific threshold for a particular indication. For other genetic sources, such as the Open Targets Genetics (OTG), Online Mendelian Inheritance in Man (OMIM), etc. associations where included based on thresholds used in previous work \cite{minikel_refining_2024}.

\begin{figure}[t!]
    \centering
    \makebox[\textwidth]{
    \begin{subfigure}{.7\textwidth}
        \centering
        \includegraphics[width=\linewidth]{figures/indirect-support/real/pigean_combined.threshold_plot.png}
        \caption{}
        \label{fig:combined-support-rs}
    \end{subfigure}
    \hspace{-10mm}
    \begin{subfigure}{.4\textwidth}
        \centering
        \includegraphics[width=\linewidth]{figures/indirect-support/real/pigean_combined.regression_analysis.png}
        \caption{}
        \label{fig:combined-support-rs-regression}
    \end{subfigure}
    }
    \makebox[\textwidth]{
        \begin{subfigure}{\textwidth}
        \centering
        \includegraphics[width=\linewidth]{figures/indirect-support/real/reproducibility.performance_by_source.png}
        \caption{}
        \label{fig:rs-comparision}
    \end{subfigure}
    }
    
    \caption{tbd}
    \label{fig:enter-label}
\end{figure}

We studied the effect of combined PIGEAN support as a function of relative success over four trait subsets. These subsets included \getvar{var__NUM_UNIQUE_PORTAL_TO_MESH_MAPPED_TRAITS} mapped indications from our A2F web portal that included full summary statistics, \getvar{var__NUM_UNIQUE_GCAT_MAPPED_TRAITS} mapped indications from the GWAS Catalog with full summary statistics, \getvar{var__NUM_UNIQUE_ORPHANET_MAPPED_TRAITS} mapped indications to associated Orphanet diseases leveraing disease-relevant gene lists, and 100 mapped indications from OMIM for direct comparision of RS on a subset of traits analyzed from previous work. Figure~\ref{fig:combined-support-rs} shows the effect of combined support as a function of RS over all clinical areas for each of these subsets. Regressing combined support over all subsets shows a strong postive correlation (Figure~\ref{fig:combined-support-rs-regression}) over both rare and common diseases ($R^2  = \getvar{var__PIGEAN_COMBINED_REGRESSION_R2}, p=\getvar{var__PIGEAN_COMBINED_REGRESSION_P_VALUE}$) that begins to fall off as combined support approaches 8, corresponding to $\sim 99\%$ that a gene is relevant for a trait. Comparing PIGEAN with a combined threshold of 8 with other sources of genetic evidence (Figure~\ref{fig:rs-comparision}), we see that PIGEAN has comparable RS for rare diseases (PIGEAN OMIM vs. OMIM) and increases RS by factor of 1.6x ($Z=\getvar{var__OTG_PIGEAN_COMBINED_Z_SCORE}$, $p=$ \getvar{var__OTG_PIGEAN_COMBINED_P_VALUE}) for common diseases over comparable subsets (e.g., GWAS Catalog). Futher, PIGEAN prioritiztions for common traits in our A2F web portal have comparable RS to traits in GWAS catalog. 

Merging trait subsets into common (i.e., GWAS Catalog and A2F portal) and rare (i.e., Orphanet and OMIM) phenotypes, we find direct and indirect can independently aid in prioritizing targets with greater likelihood of clinical success. Using direct support alone we observed a signifcant postive correlation for common diseases ($R^2=\getvar{var__PIGEAN_DIRECT_MERGED_COMMON_REGRESSION_R2}, p=$ \getvar{var__PIGEAN_DIRECT_MERGED_COMMON_REGRESSION_P_VALUE}) which confirms previous reports finding of the role of genetic support in drug target success. Since direct support for rare diseases are fixed given that genetic support is derived from known relevant gene sets, we don't consider direct support only priorization for rare diseases. Most noteably, indirect support also has a strong positve correlation with relative success for common ($R^2 = \getvar{var__PIGEAN_INDIRECT_MERGED_COMMON_REGRESSION_R2}, p=$ \getvar{var__PIGEAN_INDIRECT_MERGED_COMMON_REGRESSION_P_VALUE}) common diseases.

% Indirect support alone vs. RS (with/without genes that have high direct).

% Indirect support on programs with no genetic support. (See abstract)

% Decomposing over clinical area (area heterogenity)
%% Signal that goes from gene set betas to RS in traits that have weired indirect support. 

\subsection{Indirect Support Additively Enhances Clinical Success}
\textbf{Goal:} Argue for the unique value of indirect evidence in prioritizing clinically successful drug targets.

Outline:
\begin{itemize}
    \item Motivation
    \begin{itemize}
        \item Traditional drug target prioritization heavily favors indirect genetic evidence, and now we can support that directly with genetics
        \item Show how some successful targets lack strong direct support, but can get boosted with indirect through a pathway that a drug designer may have previously used alone.
    \end{itemize}
    \item Approach
    \begin{itemize}
        \item Use indirect-only PIGEAN score across curated drug target lists
        \item Measure clinical success enrichment (e.g., approved vs. failed targets)
    \end{itemize}
    \item Key findings
    \begin{itemize}
        \item Targets with high indirect scores show a statistically significant increase in approval likelihood (placeholder: enrichment p = XXXX, odds ratio = XXXX)
        \item No overlap with input annotations derived from clinical success data → no circularity
    \end{itemize}
    \item Implications
    \begin{itemize}
        \item Indirect support provides orthogonal, non-redundant signal
        \item Enables early-stage prioritization when direct signals are weak or absent
    \end{itemize}
\end{itemize}

\subsection{Exploration of the PIGEAN PGAG}

\begin{itemize}
    \item Ask Katy to help us visualize it
    \item Run GNN to predict
    \begin{itemize}
        \item Can we predict drug target success based just on the pheWAS + gene sets linked to it (neighbors of the gene)
        \begin{itemize}
            \item Take 8 diabetes drugs, 8 failed. just show PheWAS plus top mouse. See if can guess which ones succeed
        \end{itemize}
        \item Can we predict what rare disease the gene is involved in just based on common trait PheWAS? What if we add in pathways (gets at degree of shared biology between rare and common diseases)
        \begin{itemize}
            \item How many genes can agent guess multiple choice rare disease given PheWAS for common? PheWAS + top two gene sets ?
        \end{itemize}
        \item Describe clearly three types of gene/pheno links: direct, indirect, "GNN inferred" -- how does RS vary across these three?
        \item Can we predict effector genes better by looking at neighborhood across all traits?
    \end{itemize}
\end{itemize}

\section{Discussion}

\paragraph{Limitations}
\begin{itemize}
    \item Gene Set Biases
    \begin{itemize}
        \item Circularity from reusing GWAS-informed annotations
        \item Overrepresentation of well-studied pathways or tissues
    \end{itemize}
    \item Biological Gaps
    \begin{itemize}
        \item Underperformance in understudied traits with few known genes
    \end{itemize}    
    \item Inference Limits
    \begin{itemize}
        \item Spike-and-slab regularization may under-estimate negative effects
        \item Dependence on annotation quality; false positives if poor annotations included
    \end{itemize}
\end{itemize}

\paragraph{Conclusion}
\begin{itemize}
    \item PIGEAN introduces a principled Bayesian approach to combine direct and indirect evidence for gene prioritization and can be directly applied to drug target prioritization. 
    \item Indirect support alone is predictive of clinical success and robust to circularity.
    \item The public resource and agent tools open doors for interpretability and hypothesis generation at scale.
    \item Future work will continue to refine annotation selection, model sparsity, and dynamic integration with new data modalities.
\end{itemize}

\newpage

\bibliography{references}% common bib file
%% if required, the content of .bbl file can be included here once bbl is generated
%%\input sn-article.bbl

\newpage

%%%%% End of Main Text
 
\section{Methods}\label{sec:methods}

\subsection{PIGEAN Model}
\hypertarget{pigeanpior}{The} PIGEAN model specifies that (a) a subset of genes are ``relevant" to a trait and (b) the probability a gene is relevant to a trait is affected (additively) by the gene’s annotations. Specifically, for an arbitrary trait $T$ and genes $g_i\ldots g_N$, let $d_i \in\{0,1\}$ indicate whether gene $i$ is involved in $T$ and let $y_i=\Pr(d_i=1)$. Given a set of (binary or real valued) annotations $X_{ij}$ (where $j=1\dots M$) for $g_i$, we then model:
\[
\log\frac{y_i}{1-y_i} = \sum_j X_{ij}\beta_j + \alpha_i + \log\frac{\pi}{1-\pi} 
\]
\begin{align*}
\beta_j \sim \delta_j\mathbf{N}(0, \sigma_j^2) \quad
\alpha_i = \alpha\sum_j X_{ij} \quad
\sigma_j^2 = \sigma^2\text{Var}_i(X_{ij})^w \quad
\delta_j \sim \text{Bern}(p) \quad
\end{align*}
where $\pi$ is the prior probability that a gene is involved in disease, which we typically set to 0.05 ($\sim 1000$ genes) for complex traits and 0.01 ($\sim 200$ genes) for rare diseases, and the four parameters ($\alpha$, $\sigma$, $w$, and $p$) enforce regularization: $\sigma$ and $w$ penalize high variance (i.e. “large”) genes and annotations, $\sigma$ shrinks gene annotation effect sizes, and $p$ enforces sparsity of gene annotation effect sizes. As described in the Online Methods, we learn $p$ from the input data for trait $T$, set $w$ and $\alpha$ to limit dependence of estimands on the variance of annotations or genes, and set $\sigma$ to an empirically determined large trait-independent value (learning trait-specific $\sigma$ values is possible but, empirically, has little impact on results and introduces the potential for less robustness during model application across many traits).

We do not know $d_i$ or $\beta_j$ (they are latent variables), but we can jointly infer them conditional upon observed genetic association statistics $\mathcal{S}$. By employing Gibbs sampling, we can sample from this distribution $\prob(d, \beta \mid \mathcal{S})$ if we can sample from $\prob(d \mid \beta, \mathcal{S})$ and $\prob(\beta \mid d, \mathcal{S})=\prob(\beta \mid d)$. We can apply previous ideas[] to sample from $\prob(\beta \mid d)$ (Methods), and we consider four types of genetic association statistics of increasing complexity for calculating the conditional probability $\prob(d \mid \beta, \mathcal{S})$: rare diseases in which a subset of disease-relevant genes are "known", WES gene-level burden test statistics, genome-wide SNP summary statistics from a GWAS, and a list of only significant SNPs from a GWAS. More complicated scenarios including multiple types of summary statistics (e.g. GWAS summary statistics alongside WES burden tests of ultra-rare variants) can be handled in a straightforward manner, provided the summary statistics are independent (see Methods).

In the rare disease scenario, we assume that we are given $\prob(d \mid \mathcal{S})$, either fixed to a high value for a small number of "known" genes or estimated from qualitative confidence "tiers"[].  It is straightforward (Methods) to show that the log of the conditional odds of $d_i$ is the sum of its prior odds conditional upon $\beta$ and the unconditional gene bayes factor (which can be calculated from $\prob(d \mid \mathcal{S})$ and $\pi$), under the assumption that gene annotations included in the model are not influenced by knowledge of disease genes for $T$ (we investigate the feasibility of this assumption and the consequences of its violation below). In the WES scenario, the calculations are identical except, rather than being given $\prob(d \mid \mathcal{S})$ and converting it to gene bayes factors, we estimate the gene bayes factors from the burden effect sizes and standard errors using an approach previously developed for SNPs (Methods).

In the genome-wide GWAS summary statistics scenario, we calculate $\prob(d \mid \beta, \mathcal{S})$ using an extension of the HuGE calculator, a previously proposed set of simple rules to estimate gene-level probabilities of disease relevance conditional upon nearby SNP associations (Methods). The extensions we use within PIGEAN (referred to as HuGEX; see Methods) address some of the obvious limitations of the HuGE calculator while preserving its advantages (simplicity, speed, and interpretability). In the  scenario of partial GWAS summary statistics, we apply the HuGEX model to all genes nearby SNPs with summary statistics available, and set the remaining genes to have $\prob(d \mid \beta, \mathcal{S})=\prob(d \mid \beta)$. 

The joint posterior distribution $\Pr(d,\beta|\mathcal{S})$ inferred by PIGEAN can be distilled into four sets of summary statistics (Methods): $E[\beta_j]$, the effect size for the $j^{th}$  gene annotation; $ABF^I_i$, the “indirect” approximate Bayes Factor (ABF) that represents evidence from the annotations for gene $g_i$; $ABF^D_i$, the “direct” ABF that represents evidence from genetic associations for gene $g_i$; and $ABF_i$, the combined $ABF$ for gene $g_i$. The overall probability of disease relevance for a gene can then be expressed straightforwardly as a function of $ABF_i$ and the prior $\pi$. 

These equations offer three interpretations of the PIGEAN model (Methods). In the "prinicpled priors" interpretation, the posterior probability of a gene's disease relevance is estimated using genetic associations as data ($ABF_i^D$) and the annotations to adjust the prior (a combination of $ABF_i^I$ and $\pi$), offering bayesian approach to use both genetic associations and gene annotations to prioritize genes. In the "bottom-up" interpretation, PIGEAN generalizes the popular PoPS approach; it can be shown (Methods) that if MAGMA Z-scores are used to estimate the log-odds of $ABF_i^D$, sparsity constraints on the $\beta_j$ values are replaced by L2 regularization, and full Gibbs sampling from the posterior is replaced by a single regression, then PIGEAN reduces to PoPS --- these alterations elucidate the quantity PoPS  estimates as well as the approximations in inference and data that it makes. Finally, the "indirect genetic support" interpretation of PIGEAN demonstrates that genetic support is due not only to genetic associations nearby a gene ($ABF_i^D$) but also to genetic associations nearby other genes ($ABF_i^I$) that indirectly support the gene through shared annotations.

We implemented PIGEAN as a Python program and applied it to XXX full summary statistic GWAS datasets from the Association to Function Knowledge Portal, lists of reported significant GWAS associations for X traits from the GWAS catalog, X WES burden statistic datasets from XXX, and XXX rare disease gene lists from Orphanet (\textbf{Supplementary Table 1}). For validation purposes, we also analyzed 20 simulated “null” GWAS in which all SNPs had zero true effect sizes, and 10 simulated GWAS in which SNPs had true effect sizes consistent with the PIGEAN model (Methods). We considered a total of 144,627 gene annotations, including (Supplementary Table 2) knockout mouse phenotypes (11,502 Mammalian Phenotype Ontology terms), membership in gene sets and pathways (29,003 gene sets from the Molecular Signatures Database), occurrence in a PubMedCentral manuscript figure (19,391 gene sets from the Pathway Figure OCR project), co-occurrence with terms in PubMed abstracts (25,126 MeSH terms), network interaction partners (37,156 gene neighbors from STRING), and single-cell gene expression patterns (41,840 experiment clusters curated as part of the PoPS method).

% \subsection{PIGEAN Model Overview}
% \begin{figure}[t]
%     \centering
%     \includegraphics[width=0.5\linewidth]{figures/PIGEAN-Plates.png}
%     \caption{The full PIGEAN model.}
%     \label{fig:pigean-model}
% \end{figure}

% Our method, PIGEAN (\textbf{P}riors \textbf{I}nferred from \textbf{GE}ne \textbf{AN}notations), integrates direct genetic evidence with gene annotations to infer the relevance of genes to a given trait. This approach mirrors how a researcher might weigh both prior knowledge and direct genomic evidence when evaluating gene-disease relevance (Figure~\ref{fig:pigean-concept}). However, unlike a biomedical researcher, PIGEAN unifies these direct and indirect components under a robust probabilistic framework, providing an interpretable method for prioritizing genes for further investigation. In this section, we will briefly introduce these concepts and their formulation within the PIGEAN model (see Methods for a complete treatment). 

% \paragraph{Direct Genetic Support}
% For a given trait $T$ and genes $g_1, \dots, g_N$, let $d_i \in \{0,1\}$ indicate whether gene, $g_i$, is relevant to a trait ($d_i=1)$ or not ($d_i=0$). PIGEAN is capable of integrating direct genetic support from a variety of data sources including summary statistics from genetic association studies, WES burden tests, or even from curated disease genes or confidence tiers. In all cases, we will refer to direct genetic support ($\mathcal{S})$ as providing evidence for $d_i$, denoted $\text{Pr}(d_i \mid \mathcal{S})$. 

% We quantify this evidence using a direct approximate Bayes factor (ABF), denoted $ABF^D_i$, which represents the probability of $d_i=1$ versus $d_i=0$ based solely on genetic associations. Depending on the data source of $\mathcal{S}$, $ABG^D_i$ is calculated as follows:
% \begin{itemize}
%     \item \textbf{GWAS summary statistics:} $ABF_I^D$ is calculated using an extension of the HuGE calculator [], which calculates SNP-level associations into gene-level probabilities (referred to as HuGEX?).
%     \item \textbf{WES burden tests:} $ABF_i^D$ is estimated using effect sizes and standard errors from gene-level burden tests. 
%     \item \textbf{Rare disease gene lists:} $ABF_i^D$ is derived from prior probabilities of gene relevance, informed by known disease genes or confidence tiers. 
% \end{itemize}

% From these direct approximate Bayes Factors, a posterior probability of gene relevance to a trait can be calculated by combining these Bayes Factors with an appropriate prior probability, denoted $\hat{\pi}$, of trait relevance to yield (see Methods for complete derivation):
% \begin{equation} \label{eq:direct_BF}
%     P(d_i \mid \mathcal{S}, \hat{\pi}) = \frac{ABF_i^D \cdot \hat{\pi}}{ABF_i^D\cdot \hat{\pi} + (1-\hat{\pi})}
% \end{equation}

% where it might be a reasonable assumption to set $\hat{\pi} = 0.05$ for complex traits corresponding to $\sim 1000$ genes being relevant and $\hat{\pi}=0.01$ for rare disease corresponding to $\sim 200$ genes being relavent for a rare disease. The key insight of PIGEAN, is that it is possible to use ``indirect" knowledge from gene annotations to yield a more informed prior probablity, $\hat{\pi}$ in Equation \ref{eq:direct_BF}. 

% \paragraph{Indirect Support through Shared Gene Annotations}
% For a given gene, $g_i$, let $\mathbf{X}_i$ denote the set of all available annotations for $g_i$. These annotations can be viewed as gene sets curated from pathways databases, RNAseq enrichments, or any other modality in which genes share some function or association as determined by experiment or analysis. In this way, we can view $\mathbf{X}_i$ as a proxy for all the prior knowledge we know about or wish to consider for each gene. Next, let $X_{ij} \in \mathbf{X}_i$ represent the $j$-th annotation of gene $g_i$, such that $X_{ij}$ may be binary (inclusion or exclusion from a set) or real-valued. PIGEAN uses these annotations to model the prior probability $\hat{\pi}_i = \text{Pr}(d_i=1)$ in Equation ~\ref{eq:direct_BF} and to yield the indirect approximate Bayes Factor, denoted $ABF_i^I$, capturing the relevance of $g_i$ to a trait based on its associated annotations. Specifically:

% \begin{equation} \label{eq:indirect_prior}
%     \log \frac{\hat{\pi}_i}{1-\hat{\pi}_i} = \sum_j X_{ij}\mathbb{E}[\beta_j \mid \mathbf{d}] + \alpha_i + \log\frac{\pi}{1-\pi}
% \end{equation}
% such that $\beta_j$ represents the effect size of the $j$-th annotation and whose prior probability, $P(\beta_j)$, is given by the following spike and slab [] distribution:
% \begin{align*}
%     \beta_j \sim \delta_j\mathbf{N}(0, \sigma_j^2) \quad
%     \alpha_i = \alpha\sum_j X_{ij} \quad
%     \sigma_j^2 = \sigma^2\text{Var}_i(|X_{ij}|)^w \quad
%     \delta_j \sim \text{Bern}(p) \quad
% \end{align*}

% where $\pi$ is the background prior probability, set following the same assumptions in the direct support case: $\pi=0.05$ for complex traits and $\pi = 0.01$ for rare diseases. However, unlike the direct support case alone, the informed prior $\hat{\pi}$ can be dynamically updated based on genetic evidence from other genes that share functional annotations, by propogating through the expected annotation effect sizes, $\mathbb{E}[\beta_j \mid \mathbf{d}]$. The parameters $\alpha$, $\sigma$, $\omega$, and $p$ enforce various modes of regularization (see Methods). Specifically, $\delta_j \sim \text{Bern}(p)$ enforces gene set sparsity and $\sigma_j^2 = \sigma^2\text{Var}_i(|X_{ij}|)^\omega$ penalizes larger gene sets. 

% It is critical to note the coupling between the two latent variables $\beta_j$ and $d_i$. In Equation~\ref{eq:indirect_prior}, we need an expectation for $\beta_j$ to calculate $\hat{\pi}$ and thereby $P(d_i \mid S, \hat{\pi})$ in Equation~\ref{eq:direct_BF}, but in order to sample from $P(\beta_j \mid \mathbf{d})$, we require an estimate for each gene's relavence, $d_i$. To circumvent this circularity, we employ an iterative Gibb's sampling approach.

% \paragraph{Jointly Inferring Direct and Indirect Support}
% To jointly infer direct and indirect support, we need to sample from the joint distribution $P(\mathbf{d}, \mathbf{\hat{\pi}} \mid \mathcal{S})$. Such a sampling is possible if we can sample from the $P(\mathbf{d} \mid \beta, \mathcal{S})$ and $P(\beta \mid \mathbf{d}, \mathcal{S}) = P(\beta \mid \mathbf{d})$ since $\beta$ is the only random variable in Equation~\ref{eq:indirect_prior} and all information of $\mathcal{S}$ is encoded in $\mathbf{d}$. Therefore, given a estimate for every gene's relevance, $\mathbf{d}$, we can estimate $P(\beta_j \mid \mathbf{d})$ by logistically regressing each annotation $X_{\cdot j}$ vector on this estimate of $\mathbf{d}$ with our prior distribution, $P(\beta)$ previously specified. Using this new estimate $\tilde{\beta_j}$ of the effect size, we can use Equation~\ref{eq:direct_BF} and~\ref{eq:indirect_prior} to infer an update for overall gene relevance, $\mathbf{\tilde{d}}$. With these update equations specified for $P(\beta \mid \mathbf{d})$ and $P(\mathbf{d} \mid \beta, \mathcal{S})$, we can employ Gibbs sampling to run these updates until convergence yielding the true joint distribution. The full PIGEAN model as a directed graph is depicted in Figure~\ref{fig:pigean-model}

\subsection{Validation on Simulations}
As a negative control, we first applied PIGEAN to the 20 “null” simulated GWAS and found that, as expected, inferred gene annotation effect sizes were all close to zero (Extended Data Figure~\ref{fig:null_sim}a) as were indirect, direct, and combined gene ABFs (Supplementary Figure \ref{fig:null_sim}bcd). Similarly, when we applied PIGEAN to the 50 real GWAS but randomly permuted genes within the 166,868 gene annotations, we also observed annotation effect sizes close to zero (Supplementary  Figure \ref{fig:null_sim}e). Across our 20 null simulations and 166,868 annotations only one gene had combined gene ABF above 3 (Supplementary Figure \ref{fig:null_sim}g), with the dominant contribution coming from direct ABFs (Supplementary Figure \ref{fig:null_sim}h), and across the 50 real GWAS with permuted annotations, only one annotation had effect size above 0.004 (Supplementary Figure \ref{fig:null_sim}f). Based on these results, annotation effect sizes of 0.005 and gene combined ABFs of 3 roughly correspond to “study-wide significant” p-value thresholds (Supplementary Table 3).

We also used simulations with non-zero annotation effect sizes to understand the behavior of PIGEAN in response to different model parameters. Accuracy of PIGEAN is affected by noise from two sources: the incomplete ability of annotations to predict disease-relevant genes (“annotation noise”), and the incomplete ability of summary statistics to identify disease-relevant genes (“association noise”). Focusing first on annotation noise, we ran PIGEAN with a truncated model in which we treated each gene’s probability of disease relevance ($d_i$) as observed rather than inferred from summary statistics. These simulations yielded four main findings. First, PIGEAN in general conservatively estimates both annotation effect sizes (Supplementary Figure \ref{fig:non_null_sim}a) and gene ABFs (Supplementary Figure \ref{fig:non_null_sim}b), due to its “spike and slab” regularization that sets estimated effect sizes to zero and shrinks estimated effect sizes toward zero. Second, annotations with negative effects (i.e. that reduce genetic support) were shrunk substantially more than those with positive effects, as (consequently) were genes with ABFs < 1, an asymmetry that naturally results from the higher standard errors for negative gene set effect sizes when there are only a small number of disease-relevant genes (Supplementary Figure \ref{fig:non_null_sim}cd). Third, PIGEAN’s accuracy was also affected by simulation parameters, with more accurate inferences possible for larger simulated effect size magnitudes (larger $\sigma$ Supplementary Figure \ref{fig:sim_params}a) and (to a lesser extent) more non-zero simulated effect sizes (larger $p$; Supplementary Figure \ref{fig:sim_params}b). Additional noise added to the simulated gene ABFs affected accuracy of the inferred ABFs but (perhaps non-intuitively) did not affect accuracy of the inferred annotation effect sizes (Supplementary Figure \ref{fig:sim_params}c). Finally, due to the nature of Gibbs sampling, estimates produced by PIGEAN varied across runs, but with sufficient numbers of iterations and chains estimates were stable (Supplementary Figure \ref{fig:sim_stability}).

These simulations confirmed the importance PIGEAN’s joint inference and regularization features. First, its full joint inference procedure (in which gene and gene set values are iteratively re-estimated) increased accuracy compared to a model in which gene set effect sizes were only estimated once (Supplementary Figure \ref{fig:sim_onestep}). Second, even though simulated gene ABFs had no dependence on the number of annotations for the gene, regularization (via its  parameter) to penalize genes with many annotations was necessary to avoid biases and ABF accuracy reductions induced (as described in Supplementary Figure \ref{fig:non_null_sim}cd) by asymmetric annotation standard errors (Supplementary Figure \ref{fig:sim_induced}). Similarly, even though simulated annotation effect sizes had no dependence on the number of genes with the annotation, regularization (via its w parameter) to penalize “large” annotations was necessary to avoid biases induced by the smaller standard errors of larger annotations (Supplementary Figure \ref{fig:sim_reg}; Methods).

For simulations with a small number of non-zero annotation effects, PIGEAN was able to infer reasonably accurate estimates of its sparsity parameter ($p$) (Supplementary Figure \ref{fig:sim_param_infer}), although it significantly underestimated $p$ as the fraction of non-zero effects increased. Our attempts to estimate PIGEAN’s shrinkage parameter ($\sigma$) for individual simulations suggested that there are typically too few annotations  to estimate this parameter accurately, and we therefore tuned this parameter empirically (Methods). Fortunately, varying these parameters from their simulated values by even an order of magnitude did not dramatically reduce PIGEAN’s estimation accuracy (Supplementary Figures \ref{fig:sim_vary_sigma} and \ref{fig:sim_vary_p}).

In addition to these simulations focused exclusively on “annotation noise”, we also evaluated the impact of  “association noise” on PIGEAN’s performance – that is, the reality that (due to finite sample sizes and the vagaries of evolution) many disease-relevant genes will not lie near detectable association signals. We studied association noise in the context of GWAS summary statistics, given the well-documented difficulties in inferring effector genes from these data. For each of three annotation noise settings (Methods), we simulated five GWAS of increasing sample sizes (using linkage disequilibrium from the 1000 Genomes Project, SNP to gene linkages from a previously published method, and SNP effect size parameters from previously published simulations; Methods). As expected, the use of GWAS data rather than observed disease probabilities degraded the accuracy of PIGEAN (Supplementary Figure \ref{fig:sim_gwas_gene_set}a). With perfect knowledge of SNP to gene linkages Supplementary Figure \ref{fig:sim_gwas_gene_set}b), less annotation noise Supplementary Figure \ref{fig:sim_gwas_gene_set}c), or larger sample sizes Supplementary Figure \ref{fig:sim_gwas_gene_set}d), the accuracy of PIGEAN increased, but only to limited extent. Accuracy was much higher for marginal annotation effect sizes than for joint effect sizes Supplementary Figure \ref{fig:sim_gwas_gene_set}e), which is a common feature of multivariate models with correlated predictors. When PIGEAN inferred its regularization parameters, rather than being given the true parameters used in simulations, accuracy degraded further Supplementary Figure \ref{fig:sim_gwas_gene_set}f). Notably, however, in all scenarios PIGEAN conservatively estimated annotation effect sizes, suggesting that the overwhelming impact of GWAS annotation noise is to decrease power rather than to increase false positives.

We conclude that, under idealized simulations, PIGEAN behaves broadly as expected. Three noteworthy properties are (a) the difficulty of inferring negative genetic support (due to known properties of GWAS but perhaps unappreciated properties of models with few disease-relevant genes); (b) the importance of regularization and iterative inference in gene set analysis, and (c) it conservative nature when confronted with annotation noise and and association noise.

\subsection{Validation on Real Data}
\paragraph{Selection of Validation Phenotypes}
%TODO: Also include bottom line and/or credible set information for each phenotype to prove we looked at phenos with varying GWAS significance levels
For common traits, we selected 13 phenotypes (Table~\ref{tab:common-validation-traits}) encompassing both continuous and dichotomous traits with GWAS data of varying significance but which we had curated gold-standard genesets for validation. This diversity allows us to evaluate PIGEAN and respective PIGEAN models under a broad range of real-world conditions. For rare traits, we chose 40 phenotypes (Table~\ref{tab:rare-validation-traits}) that span multiple organ systems, inheritance patterns, and pathophysiological mechanisms. This set includes both specific rare diseases (with more than 10 associated genes) and groups of disorders representing higher-level disease clusters. To ensure a balanced representation without excessive overlap in gene lists, we verified that the maximum Jaccard similarity among groups of disorders is 0.14 and 0.11 among specific diseases, with average similarities of 0.00653 and 0.0016, respectively. By covering a wide range of disease etiologies and genetic architectures, we can thoroughly assess our method's performance across diverse traits.

\begin{table}[ht]
\centering
\footnotesize
\begin{tabular}{l c c}
\hline
\textbf{Disease Name} & \textbf{\# Sig GWAS Loci} & \textbf{\# Credible Sets} \\
\hline
\href{https://a2f.hugeamp.org/phenotype.html?phenotype=AF}{Atrial Fibrillation} & TBD & TBD \\
\href{https://a2f.hugeamp.org/phenotype.html?phenotype=ALLOA}{Osteoarthritis} & TBD & TBD \\
\href{https://a2f.hugeamp.org/phenotype.html?phenotype=AnyCVD}{Cardiovascular Disease} & TBD & TBD \\
\href{https://a2f.hugeamp.org/phenotype.html?phenotype=Eczema}{Eczema} & TBD & TBD \\
\href{https://a2f.hugeamp.org/phenotype.html?phenotype=HDL}{High-Density Lipoprotein Cholesterol (HDL)} & TBD & TBD \\
\href{https://a2f.hugeamp.org/phenotype.html?phenotype=HF}{Heart Failure} & TBD & TBD \\
\href{https://a2f.hugeamp.org/phenotype.html?phenotype=HYPERTENSION}{Hypertension} & TBD & TBD \\
\href{https://a2f.hugeamp.org/phenotype.html?phenotype=LDL}{Low-Density Lipoprotein Cholesterol (LDL)} & TBD & TBD \\
\href{https://a2f.hugeamp.org/phenotype.html?phenotype=MDD}{Major Depressive Disorder} & TBD & TBD \\
\href{https://a2f.hugeamp.org/phenotype.html?phenotype=MI}{Myocardial Infarction} & TBD & TBD \\
\href{https://a2f.hugeamp.org/phenotype.html?phenotype=RhA}{Rheumatoid Arthritis} & TBD & TBD \\
\href{https://a2f.hugeamp.org/phenotype.html?phenotype=SCZ}{Schizophrenia} & TBD & TBD \\
\href{https://a2f.hugeamp.org/phenotype.html?phenotype=T2D}{Type 2 Diabetes} & TBD & TBD \\
\hline
\end{tabular}
\caption{Common Disease Validations with hyperlinked names and placeholders for GWAS loci and credible sets.}
\label{tab:common-validation-traits}
\end{table}


\begin{table}[ht]
\centering
\footnotesize
\begin{tabular}{l l l l c}
\hline
\textbf{Disease Name} & \textbf{Orphanet ID} & \textbf{Type} & \textbf{\# Genes} \\
\hline
Neuromuscular disease & \href{https://www.orpha.net/en/disease/detail/68381}{68381} & Group & 374 \\
Inherited arrhythmogenic cardiomyopathy & \href{https://www.orpha.net/en/disease/detail/247}{247} & Group & 13 \\
Primary immunodeficiency & \href{https://www.orpha.net/en/disease/detail/101997}{101997} & Group & 328 \\
Mucopolysaccharidosis & \href{https://www.orpha.net/en/disease/detail/79213}{79213} & Group & 11 \\
Sphingolipidosis & \href{https://www.orpha.net/en/disease/detail/79225}{79225} & Group & 17 \\
Ciliopathy & \href{https://www.orpha.net/en/disease/detail/363250}{363250} & Group & 282 \\
Rare sleep disorder & \href{https://www.orpha.net/en/disease/detail/68354}{68354} & Group & 14 \\
Rare pervasive developmental disorder & \href{https://www.orpha.net/en/disease/detail/168778}{168778} & Group & 25 \\
Sexchromosome anomaly & \href{https://www.orpha.net/en/disease/detail/98155}{98155} & Group & 21 \\
Autosomal trisomy & \href{https://www.orpha.net/en/disease/detail/98130}{98130} & Group & 16 \\
Hereditary ataxia & \href{https://www.orpha.net/en/disease/detail/183518}{183518} & Group & 182 \\
Peroxisomal disease & \href{https://www.orpha.net/en/disease/detail/68373}{68373} & Group & 28 \\
Mitochondrial disease & \href{https://www.orpha.net/en/disease/detail/68380}{68380} & Group & 208 \\
Disorder of purine or pyrimidine metabolism & \href{https://www.orpha.net/en/disease/detail/79224}{79224} & Group & 50 \\
Congenital myasthenic syndrome & \href{https://www.orpha.net/en/disease/detail/590}{590} & Group & 27 \\
Rare hypertrophic cardiomyopathy & \href{https://www.orpha.net/en/disease/detail/217569}{217569} & Group & 81 \\
Amyloidosis &\href{https://www.orpha.net/en/disease/detail/69}{69} & Group & 11 \\
Hemoglobinopathy & \href{https://www.orpha.net/en/disease/detail/68364}{68364} & Group & 12 \\
Spondyloepiphyseal dysplasia & \href{https://www.orpha.net/en/disease/detail/253}{253} & Group & 38 \\
Ectodermal dysplasia syndrome & \href{https://www.orpha.net/en/disease/detail/79373}{79373} & Group & 99 \\
Amyotrophic lateral sclerosis & \href{https://www.orpha.net/en/disease/detail/803}{803} & Specific & 37 \\
Zellweger syndrome & \href{https://www.orpha.net/en/disease/detail/912}{912} & Specific & 13 \\
Fanconi anemia & \href{https://www.orpha.net/en/disease/detail/84}{84} & Specific & 22 \\
RomanoWard syndrome & \href{https://www.orpha.net/en/disease/detail/101016}{101016} & Specific & 19 \\
Cystic fibrosis & \href{https://www.orpha.net/en/disease/detail/586}{586} & Specific & 19 \\
Noonan syndrome & \href{https://www.orpha.net/en/disease/detail/648}{648} & Specific & 14 \\
Acute promyelocytic leukemia & \href{https://www.orpha.net/en/disease/detail/520}{520} & Specific & 13 \\
Bulbospinal muscular atrophy & \href{https://www.orpha.net/en/disease/detail/206701}{206701} & Specific & 16 \\
Glioblastoma & \href{https://www.orpha.net/en/disease/detail/360}{360} & Specific & 12 \\
Williams syndrome & \href{https://www.orpha.net/en/disease/detail/904}{904} & Specific & 20 \\
Kallmann syndrome  & \href{https://www.orpha.net/en/disease/detail/478}{478} & Specific & 21 \\
Hirschsprung disease & \href{https://www.orpha.net/en/disease/detail/388}{388} & Specific & 14 \\
Brugada syndrome  & \href{https://www.orpha.net/en/disease/detail/130}{130} & Specific & 22 \\
Dyskeratosis congenita & \href{https://www.orpha.net/en/disease/detail/1775}{1775} & Specific & 13 \\
Behet disease & \href{https://www.orpha.net/en/disease/detail/117}{117} & Specific & 15 \\
Jeune syndrome & \href{https://www.orpha.net/en/disease/detail/474}{474} & Specific & 11 \\
Congenital hypothyroidism & \href{https://www.orpha.net/en/disease/detail/442}{442} & Specific & 28 \\
DiamondBlackfan anemia & \href{https://www.orpha.net/en/disease/detail/124}{124} & Specific & 26 \\
Leber hereditary optic neuropathy & \href{https://www.orpha.net/en/disease/detail/104}{104} & Specific & 12 \\
WalkerWarburg syndrome & \href{https://www.orpha.net/en/disease/detail/899}{899} & Specific & 14 \\
\hline
\end{tabular}
\caption{Validation sets containing 40 diseases (20 groups of disorders and 20 specific disorders), with Orphanet links and placeholders for gene counts.}
\label{tab:rare-validation-traits}
\end{table}

\subsection{Gene Set Library Curation}
% Section about MeSH Text Mining models
For text mining models, we curated text mining based gene sets from three libraries: Gene-Mesh Co-Occurance via Pubmed publications, String DB, and RummaGene. For the Gene-MeSH co-occurance gene sets, we downloaded Gene to PubMed article mappings from the PubTator3 \cite{wei_pubtator_2024} FTP server\footnote{The gene2pubtator3.gz file is available at the following \href{https://ftp.ncbi.nlm.nih.gov/pub/lu/PubTator3/}{link} and was downloaded on May 13, 2025 for all results presented.} which includes results are from AIONER \cite{luo_aioner_2023} for Named Entity Recognition and GNorm2 \cite{wei_gnorm2_2023} for normalization. In order to contruct the MeSH term to Pubmed co-occurances, we downloaded the full Pubmed Baseline containing all pubmed information for 2025\footnote{The Pubmed baseline XML files are available at the following \href{https://ftp.ncbi.nlm.nih.gov/pubmed/baseline/}{link} and the January 10, 2025 release was used for all results presented.}. From these set of co-occurances, we then created a master gene set library that would have included publications that were based on genetics including GWAS. Therefore, we curated a list of MeSH terms related to GWAS as seen in Table~\ref{tab:gwas-mesh-terms}. We then formed three MeSH-based gene set libraries: (1) a library in which all pubmed publications are considered, (2) a library in which all GWAS and genetics-based papers are excluded (no circularity), and (3) a library that contains only GWAS and genetics-based papers (fully circular). We then used these three libraries to study the effect of circularity explicitly. Further, we waited each gene in their associated geneset by a normalized jaccard index and which each MeSH term has an associated gene set. MeSH terms that had more than 2000 genes associated with them were excluded.

\begin{table}[h]
\centering
\begin{tabular}{ll}
\toprule
\textbf{MeSH ID} & \textbf{Concept Name} \\
\midrule
\href{https://meshb.nlm.nih.gov/record/ui?ui=D055106}{D055106} & Genome-Wide Association Study \\
\href{https://meshb.nlm.nih.gov/record/ui?ui=D056726}{D056726} & Genetic Association Studies \\
\href{https://meshb.nlm.nih.gov/record/ui?ui=D000073336}{D000073336} & Whole Genome Sequencing \\
\href{https://meshb.nlm.nih.gov/record/ui?ui=D000073359}{D000073359} & Whole Exome Sequencing \\
\href{https://meshb.nlm.nih.gov/record/ui?ui=D016022}{D016022} & Case-Control Studies \\
\href{https://meshb.nlm.nih.gov/record/ui?ui=D015894}{D015894} & Genome, Human \\
\href{https://meshb.nlm.nih.gov/record/ui?ui=D015810}{D015810} & Linkage Disequilibrium \\
\href{https://meshb.nlm.nih.gov/record/ui?ui=D020022}{D020022} & Genetic Predisposition to Disease \\
\href{https://meshb.nlm.nih.gov/record/ui?ui=D020641}{D020641} & Polymorphism, Single Nucleotide \\
\href{https://meshb.nlm.nih.gov/record/ui?ui=D005838}{D005838} & Genotype \\
\href{https://meshb.nlm.nih.gov/record/ui?ui=D005787}{D005787} & Gene Frequency \\
\href{https://meshb.nlm.nih.gov/record/ui?ui=D040641}{D040641} & Quantitative Trait Loci \\
\href{https://meshb.nlm.nih.gov/record/ui?ui=D023281}{D023281} & Genomics \\
\href{https://meshb.nlm.nih.gov/record/ui?ui=D057182}{D057182} & Mendelian Randomization Analysis \\
\href{https://meshb.nlm.nih.gov/record/ui?ui=D059014}{D059014} & High-Throughput Nucleotide Sequencing \\
\href{https://meshb.nlm.nih.gov/record/ui?ui=D056915}{D056915} & DNA Copy Number Variations \\
\href{https://meshb.nlm.nih.gov/record/ui?ui=D052707}{D052707} & Genomic Instability \\
\href{https://meshb.nlm.nih.gov/record/ui?ui=D004252}{D004252} & DNA Mutational Analysis \\
\href{https://meshb.nlm.nih.gov/record/ui?ui=D008040}{D008040} & Genetic Linkage \\
\href{https://meshb.nlm.nih.gov/record/ui?ui=D019656}{D019656} & Loss of Heterozygosity \\
\bottomrule
\end{tabular}
\caption{MeSH terms excluded to avoid literature derived from GWAS, variant, and sequencing-based genetic analyses.}
\label{tab:gwas-mesh-terms}
\end{table}

\subsection{Relative Success as a Metric}
Discuss relative success methods

\subsection{Common Disease GWAS}
Discuss GWAS Catalog curation and portal curation. Discuss mapping to MESH.

\subsection{Rare Disease Curation}
Discuss running PIGEAN on gene lists. Discuss mapping to MeSH.

\section{Extended Data} \label{sec:extended-data}
% % Deplicate this as a placeholder for figures
% \begin{figure}[h!]
%     \centering
%     \begin{tikzpicture}
%         % Draw a light gray rectangle as the placeholder
%         \fill[gray!20] (0, 0) rectangle (\textwidth, 4); 
%         % Add text in the center of the rectangle
%         \node[align=center] at (\textwidth/2, 2) {\Large Placeholder Figure};
%     \end{tikzpicture}
%     \caption{Caption describing the placeholder.}
%     \label{fig:placeholder}
% \end{figure}

%JF COMMENT: For the causal evidence graph, why is there an edge from S to X? Doesn’t that imply that the annotation depends on the association, which is when we think there IS a problem with circularity?
\begin{figure}[t!]
    \centering
    \includegraphics[width=\linewidth]{figures/pigean-causal-model-figure-1.pdf}
    \caption{(a) A causal model relating genetics data such as GWAS, WES, and rare disease-gene links with gene annotation (i.e., gene sets). The figure is decomposed over all genes, $g_1, \dots, g_i$ that have some latent biology $U_i$ relevant to a trait. We introduce a latent random variable derived from $U_i$ that directly models if a gene is relevant to a trait. Further, we have the genetic information $S_i$ and gene annotation $X_i$ for each related to each other and $U_i$ in the direct graph. (b) An analytical solution to the causal model showing that highly correlated genesets and genetic information only increases the standard error on the effect sizes of the model, but does not degrade or enhance accuracy of the model.}
    \label{fig:pigean-causal-model}
\end{figure}


% Start of real figures
\begin{figure}[t!]
    \centering
    \begin{subfigure}{\textwidth}
        \centering
        \includegraphics[width=\linewidth]{figures/indirect-support/real/all.pops_vs_pigean_comparision.summary_over_phenotype.precision.png}
        \caption{}
        \label{fig:pops-v-pigean-precision-all}
    \end{subfigure}
    \hfill
    \begin{subfigure}{\textwidth}
        \centering
        \includegraphics[width=\linewidth]{figures/indirect-support/real/all.pops_vs_pigean_comparision.summary_over_phenotype.recall.png}
        \caption{}
        \label{fig:pops-v-pigean-recall-all}
    \end{subfigure}
    \caption{PoPS vs. Pigean over large set of various annotation source (Methods) comparison using gold-standard gene sets informed by an LLM analysis using precision and recall @ $K$ over 13 complex traits.}
    \label{fig:pops-vs-pigean-goldstandard-all}
\end{figure}

\begin{figure}
    \centering
    \begin{subfigure}{\textwidth}
        \centering
        \includegraphics[width=\linewidth]{figures/indirect-support/real/depict.pops_vs_pigean_comparision.summary_over_phenotype.precision.png}
        \caption{}
        \label{fig:pops-v-pigean-precision-depict}
    \end{subfigure}
    \hfill
    \begin{subfigure}{\textwidth}
        \centering
        \includegraphics[width=\linewidth]{figures/indirect-support/real/depict.pops_vs_pigean_comparision.summary_over_phenotype.recall.png}
        \caption{}
        \label{fig:pops-v-pigean-recall-depict}
    \end{subfigure}
    \caption{PoPS vs. Pigean over depict annotations comparison using gold-standard gene sets informed by an LLM analysis using precision and recall @ $K$ over 13 complex traits.}
    \label{fig:pops-vs-pigean-goldstandard-depict}
\end{figure}

\begin{figure}
    \centering
    \includegraphics[width=1\linewidth]{figures/indirect-support/real/2__full_oa4ol__rare.gene_set_hist.pdf}
    \caption{Caption}
    \label{fig:rare-full-model-genesets}
\end{figure}

\begin{figure}
    \centering
    \includegraphics[width=\linewidth]{figures/indirect-support/real/v1_all-areas.pigean_rs_curves.png}
    \caption{Relative success of PIGEAN supported target-indication pairs overall clinical areas for common (from curated GWAS and GWAS catalog traits) and rare diseases curated from Orphanet.}
    \label{fig:rs-curves}
\end{figure}

\begin{figure}
    \centering
    \includegraphics[width=\linewidth]{figures/indirect-support/real/v1_all-areas_relative_success_heatmap.png}
    \caption{Heatmap of the relative success of PIGEAN supported target-indications as a function of both direct and indirect support overall clinical areas. }
    \label{fig:rs-heatmap}
\end{figure}

\begin{figure}
    \centering
    \includegraphics[width=\linewidth]{figures/indirect-support/real/v1_all-areas_relative_success_vs_indirect.png}
    \caption{Relative success as a function of just indirect support (Direct support $<=$ 5.)}
    \label{fig:rs-indirect}
\end{figure}

\begin{figure}
    \centering
    \includegraphics[width=\linewidth]{figures/indirect-support/real/suppfig1.png}
    \caption{}
    \label{fig:null_sim}
\end{figure}

\begin{figure}
    \centering
    \includegraphics[width=\linewidth]{figures/indirect-support/real/suppfig2.png}
    \caption{}
    \label{fig:non_null_sim}
\end{figure}

\begin{figure}
    \centering
    \includegraphics[width=\linewidth]{figures/indirect-support/real/suppfig3.png}
    \caption{}
    \label{fig:sim_params}
\end{figure}

\begin{figure}
    \centering
    \includegraphics[width=\linewidth]{figures/indirect-support/real/suppfig4.png}
    \caption{}
    \label{fig:sim_stability}
\end{figure}

\begin{figure}
    \centering
    \includegraphics[width=\linewidth]{figures/indirect-support/real/suppfig5.png}
    \caption{}
    \label{fig:sim_onestep}
\end{figure}

\begin{figure}
    \centering
    \includegraphics[width=\linewidth]{figures/indirect-support/real/suppfig6.png}
    \caption{}
    \label{fig:sim_reg}
\end{figure}

\begin{figure}
    \centering
    \includegraphics[width=\linewidth]{figures/indirect-support/real/suppfig7.png}
    \caption{}
    \label{fig:sim_induced}
\end{figure}

\begin{figure}
    \centering
    \includegraphics[width=\linewidth]{figures/indirect-support/real/suppfig8.png}
    \caption{}
    \label{fig:sim_param_infer}
\end{figure}

\begin{figure}
    \centering
    \includegraphics[width=\linewidth]{figures/indirect-support/real/suppfig9.png}
    \caption{}
    \label{fig:sim_vary_sigma}
\end{figure}

\begin{figure}
    \centering
    \includegraphics[width=\linewidth]{figures/indirect-support/real/suppfig10.png}
    \caption{}
    \label{fig:sim_vary_p}
\end{figure}

\begin{figure}
    \centering
    \includegraphics[width=\linewidth]{figures/indirect-support/real/suppfig11.png}
    \caption{}
    \label{fig:sim_gwas_gene_set}
\end{figure}

\section{Supplementary Information}
% This is where supplementary tables and things go


%%===========================================================================================%%
%% If you are submitting to one of the Nature Portfolio journals, using the eJP submission   %%
%% system, please include the references within the manuscript file itself. You may do this  %%
%% by copying the reference list from your .bbl file, paste it into the main manuscript .tex %%
%% file, and delete the associated \verb+\bibliography+ commands.                            %%
%%===========================================================================================%%

%%% PARKING LOT
\begin{comment}






\section{Methods}

Our strategy to infer mechanisms underlying GWAS associations for a trait has two steps (Figure 1). In the first step, we use PIGEAN to identify genes with genetic support, and gene sets that impact genetic support, from the GWAS summary statistics (p-values and effect sizes for each SNP). In the second step, we assemble a (sparse) matrix of genes with high genetic support and gene sets with large effects on genetic support (with entries indicating which genes are within each set), we run bayesian non-negative matrix factorization with automatic relevance determination to identify latent factors each characterized by sparse and positive loadings for each gene and each gene set, and we use the gene sets and their loadings to label each factor using a large language model (LLM). The interpretation of these results is collectively (a) estimates of genetic support for the trait for each gene, separated into direct (significant associations linked to the gene) and indirect (disease-relevant annotations of the gene) genetic support estimates; (b) a list of gene sets that affect indirect genetic support, with effects characterized both by marginal values (when the gene set is considered in isolation as in existing gene set enrichment approaches) and joint values (when all gene sets are considered together); and (c) a list of latent factors that simultaneously and quantitatively group genes and gene sets into higher order mechanisms. The latent factors are designed to capture higher order biological mechanisms for which existing gene set libraries are noisy, biased, and incomplete proxies.

There are three novel components to this approach. The first novel component is the PIGEAN model, which probabilistically generalizes the widely used PoPs approach for effector gene prioritization.   assumes that a subset of genes are involved in each trait and that the probability a gene is involved in a trait is affected (additively) by the gene’s annotations. Specifically for an arbitrary trait $T$, genes $i=1, \dots, N$, and annotations $j=1, \dots, M$, let $d_i$ indicate whether gene $i$ is involved in $T$, let $j$ indicate the degree to which annotation $j$ affects this probability, and let $X_{ij}$ indicate the value of annotation $j$ for gene $i$. We first model the log-odds of the probability di as a weighted sum of the annotations for gene $i$,

\begin{align*}
\prob(z_{ik} \mid d_i) &= f(d_i) \\
d_i &\sim \text{Bern}(y_i) \\
\log\frac{y_i}{1-y_i} &= \sum_j X_{ij}\beta_j + \alpha_i + \log\frac{\pi}{1-\pi} \\
\beta_j &\sim \delta_j\mathbf{N}(0, \sigma_j^2) \\
\alpha_i &= \alpha\sum_j X_{ij} \\
\sigma_j^2 &= \sigma^2\text{Var}_i(X_{ij})^w \\
\delta_j &\sim \text{Bern}(p) \\
\end{align*}

Each node $\beta_j$ corresponds to a gene set $j$, and
\begin{equation*}
    X_{ij} =
\begin{cases}
1 & g_i \in \text{gene set } j \\
0 & \text{otherwise}
\end{cases}
\end{equation*}

where $\pi$ the prior probability that a gene is involved in disease (we typically set this to 0.05, corresponding to roughly 1000 genes, for complex traits) and the four parameters ($\alpha$, $\sigma$ , $w$, and $p$) enforce regularization: $\sigma$ and $w$ penalize high variance (i.e. “large”) genes and annotations,  shrinks gene annotation effect sizes, and $p$ enforces sparsity of gene annotation effect sizes. We learn values for $\sigma$ and $p$ from the input data (Methods) and set $w$ and $\alpha$ to enforce independence between the variance of annotations or genes and estimands.

A third interpretation is as a statistical generalization of the popular PoPS approach, with (roughly) a dichotomous model of gene disease relevance in place of MAGMA Z-scores, sparsity constraints on the $j$ values in place of L2 regularization, and full Gibbs sampling that enables iterative re-estimation of both gene and gene set values in place of a single regression and cross-validation.

$\sigma$ is hard to learn with a small number of gene sets. Therefore, we set it to a top value of 20\% for mouse genes (from Flannick et all) and assume 1\% of them are causal. This is supported by learning p across many traits (see Methods). Future work could estimate sigma/sigma\_power from data, as the sigma power does depend on the $\sigma$ (Figure)

We do not know which genes are involved in each trait, nor do we know the degree to which each annotation affects this probability ($d_i$ and $\beta_j$ are unobserved). However, if we further assume that observed association strength between a trait and a variant is affected by whether or not the variant is linked to a gene is involved in the trait, we can jointly infer $d_i$ and $beta_j$ conditional upon observed variant association statistics $z_k$ – that is, $\prob(d_i, \beta_j \mid z_k)$. By employing Gibbs sampling, we can sample from this distribution if we analytically solve for $\prob(d_i \mid \beta_j, z_k)$ and $\prob(\beta_j \mid d_i, z_k)=\prob(\beta_j \mid d_i)$. We can apply previous ideas to calculate $\prob(\beta_j \mid d_i)$ (Methods), while the model for $\prob(d_i\mid \beta_j, z_k)$ is more complex. Ideally, this model would incorporate (a) evidence linking variants to genes; (b) the $z_k$ values for variants linked to a gene; (c) variant annotations (e.g. bioinformatic or epigenomic); and (d) LD relationships among variants. For the current version of PIGEAN, we use a simpler model for $\prob(d_i \mid \beta_j, z_k)$ that builds upon previously proposed Human Genetic Evidence (HuGE) scores for calculating genetic support from common and rare variant associations[Dornbos]. We incorporate (see Methods) several extensions to the previously common variant HuGE score model that address some of its obvious limitations while preserving its advantages (namely, simplicity, speed, and interpretability).\\

Equation for rare disease sum:
\begin{align*}
\frac{\Pr(d\mid \beta,\mathcal{S})}{\Pr(\neg{d}\mid \beta,\mathcal{S})} &= \frac{\Pr(\beta,\mathcal{S}\mid  d)\Pr(d)}{\Pr(\beta,\mathcal{S}\mid \neg 
 d)\Pr(\neg d)} \\
&= \frac{\Pr(\beta,\mathcal{S}\mid d)\Pr(d)}{\Pr(\beta,\mathcal{S} \mid \neg d)\Pr(\neg d)} \\
&= \frac{\Pr(\beta,\mathcal{S}\mid d)\Pr(d)}{\Pr(\beta,\mathcal{S} \mid \neg d)\Pr(\neg d)} \\
&= \frac{\Pr(\beta \mid d)\Pr(\mathcal{S}\mid d)\Pr(d)}{\Pr(\beta \mid \neg d)\Pr(\mathcal{S} \mid \neg d)\Pr(\neg d)} \\
\end{align*}

MAKE THIS MORE OF A FRAMEWORK FOR UNDERSTANDING top-down/bottom-up or convergence/locus etc. Intuitively, PIGEAN can be viewed from four perspectives. First, its model is analogous to that used in models of polygenic SNP effects, with gene annotations playing the role of SNPs, genes playing the role of individuals, and a binary/logistic model used in place of a continuous model for the gene “phenotype”; in fact, the method used for sampling from $(\beta_j \mid d_i,z_k)$ is identical to that used by the LD-pred algorithm for sampling from the joint distribution of SNP effects. A second interpretation of PIGEAN is as a method to calculate the implication of assumptions about causal genes given observed SNP associations and “prior” knowledge about the gene’s function; in this framing, the posterior odds of a gene’s disease-relevance can be expressed in a Bayesian framework as 
\begin{equation}
    \frac{\prob(d_i \mid \beta_j, z_k)}{\prob(\neg d_i \mid \beta_j,z_k)} =
    \frac{\prob(z_k\mid d_i)}{\prob(z_k \mid \neg d_i)} \cdot \frac{\prob(d_i \mid \beta_j)}{\prob(\neg d_i \mid \beta_j)}
\end{equation}
or an update of the gene’s prior (given its annotations) after observing data (the nearby SNP associations). Finally, PIGEAN can be viewed as calculating genetic support for a gene that statistically integrates both “direct” support from nearby variant associations and “indirect” support from variant associations nearby other genes; these associations nearby other genes can indirectly support the gene if the other genes share annotations with the gene.

Given an input GWAS and collection of gene annotations, PIGEAN infers posterior distributions over $d_i$ for all genes and $\beta_j$ for all gene annotations. We distill these into four sets of summary statistics (Methods): $E[\beta_j]$, the effect size for each gene annotation $\beta_j$; $ABF^I_i$, the “indirect” approximate Bayes Factor (ABF) that represents evidence from the annotations for gene $i$; $ABF^D_i$, the “direct” ABF that represents evidence from genetic associations for gene i; and ABFi, the combined $ABF$ for gene $i$. Approximately, 
\begin{align*}
    \frac{\prob(d_i\mid \beta_j, z_k)}{\prob(\neg d_i \mid \beta_j, z_k)} \approx \frac{\pi}{1-\pi} \cdot ABF_i &= \frac{\pi}{1-\pi} \cdot ABF^I_i \cdot ABF^D_i \\
    &= \frac{\pi}{1-\pi} \sum_j X_{ij} \cdot E[\beta_j] \cdot ABF^D_i
\end{align*}




\end{comment}

%%% NOTES
\begin{comment}

Todo validation of rare indirect support:
a) Remove genes from the input list then predict (need to remove any input gene sets correlated with the gene list above some threshold)
b) Predict additional genes for CRDC phenotypes, then see whether they are enriched for burden test results (talk to Kyuryung about this)


To understand this performance, we decomposed these results. To compare improvements in indirect support, we ran cross validation. Did as good with far fewer parameters. This was assisted by the iterative updates. But the full model was necessary to get the high performance with this much sparser model. To compare improvements with direct support, we looked at MAGMA. Enrichments were XXX. The top ranked sets were XXX. The joint effects though were more sparse. Joint effects were much better than for PoPS.

Performance as a function of parameters. This would be good just for us to know. And might be worth a supp figure in the manuscript.

Show PIGEAN can be run with parameters that reproduce PoPS
Show genes prioritized by PIGEAN JUST BY PRIORS are as good or better than those prioritized by PoPS. For this we need cross_val PIGEAN priors vs. POPS priors.out
Curated positive controls (AuROC) – Ryan has these
Significant exome associations
Cross-validation on GWAS – within PoPS/cross val: see whether MAGMA Z or huge log_bf is better predicted by PoPS or PIGEAN scores
Others?
Show gene set values predicted by PIGEAN are better than PoPS
For this, we can use the main gene sets for PIGEAN and the meain gene_set_stats for PoPS
Ask LLM which of top 50 gene sets by PIGEAN and PoPS are relevant to the trait (expert curation)
Others?
Show that the PIGEAN COMBINED scores are better than either the direct (log_bf) or indirect (priors) alone. For this we can use the main (non cross val)
Same metrics
Plus predict the future: on GWAS from 5 years ago, regress MAGMA Z-score from recent GWAS on PoPS/PIGEAN score + old MAGMA Z-score
Show mechanisms by NNMF are better for PIGEAN combined > PIGEAN priors > PoPS
Run factoring (NNMF) on PoPS and PIGEAN and show that factors are better – one way is some inherent measure of factor quality, another is ability to predict other traits

We next evaluated the gene sets inferred by PIGEAN on real data.
Figure: download PIGEAN and MAGMA across all traits in the portal. (a) Take top 10 MsigDB gene sets for each, ask LMM how many match. (b) Have Patrick add newest mouse gene sets to MAGMA and PIGEAN. Find out in Monarch how to map EFO to MPO. Then calculate how those are ranked relative to the total. (c) as a function of rank, calculate how many other traits are enriched. Metric should be how many are enriched in traits that are not correlated with the primary trait
Analysis: Pops features have to be grouped -- too correlated. Compare plot of features due to sparsity

We next compared PIGEAN to PoPS, the state-of-the-art method.
Analysis: First, set up two runs for PoPS and PIGEAN. One with all PoPS gene sets, one with just depict. Then, run this across 10 traits or so with effector gene lists available and exome sequencing results available.
Effector gene list analysis:

First verify that PIGEAN in PoPS mode is highly correlated with PoPS, and good enough to use as our PoPS implementation for subsequent tests – USE DEPICT
Download test sets from latest EGL LLM paper (nearest, coding, GWAS catalog, drug target, OpenTargets) and add in our Egls – THIS IS THE GS YOU ARE USING
Obtain prioritized effector genes. – RYAN HAS THESE – BEWARE THOUGH NOT REALLY GS
First, compare the lists. What is the Venn Diagram? Look closely at those in one but not the other, and potentially add in information to PIGEAN if needed
Then, ask how many of these genes are on the effector gene lists.
Also do all performance comparisons as in the EGL LLM paper
Compare exome association enrichments for the effector genes
Separate results out for combined vs. priors vs. priors LOCO
Global ranking analysis:
Enrichment of gold standard gene lists (separate combined vs. priors vs. priors LOCO)
Use pops expression only, predict drug targets or monogenic or mouse curated
Enrichment of exome associations (using Alex metric)
Run on older GWAS, then predict future MAGMA score controlling for current MAGMA score
Run on older GWAS, then predict future SNP p-value controlling for current SNP p-value
Cross-validation analysis:
First run LOCO for PIGEAN
Then ask how well this predicts GWAS MAGMA

We asked which novel aspects of PIGEAN were most responsible for the performance improvements.
Which aspect gives the most improvement? HuGE, sparsity, iteration, combining priors + HuGE
Betas original vs final. Two axes: first vs. after full gibbs. Uncorrected vs corrected. Plus those two axes together
Look at pops gene sets high in pops alone but not in all. Did orig beta change? Or parameters?

We also evaluated the performance on exome analyses.
Run just on exome, and use this to predict GWAS
Find exome studies, inflate standard errors of gene-level associations (and resample the betas from the uninflated distributions), run on this and predict future
Find actual older T2D WES (and the smaller UKB 50K) and predict future discoveries from that as well

Sensitivity analyses (effect of parameters):
Run these with different sigma power and see if results change
Try detect-power and non-detect power for this – gene sets look better for detect power
Try also effects of sigma, effects of huge parameters (mainly the sig_posterior – try fixed threshold 5e-8 and fixed signal percentile, and try using this for all gene support or only for signal inference)
Rerun removing any gene set with >1 gene in the gold standard set
Rerun using traits from 2020 to avoid overfitting

To extend these effector gene prioritization ideas to a more complete approach to gene and mechanism discovery for complex traits, several gaps remain. First, many (if not most) candidate genes do not lie within a GWAS locus, as do many genes “known” to be relevant to the trait. No existing method can statistically integrate association evidence for a gene outside of GWAS loci with its annotations to estimate the probability it is involved in the trait. Second, existing methods produce gene values comparable within a GWAS locus (i.e. via ranking of genes) but difficult to interpret across loci or traits. Third, existing methods are heuristic and non-trivial to statistically incorporate into downstream analyses such as combination with GWAS or WES association evidence; one approach has shown the potential for machine learning to integrate PoPS results with SNP-to-gene data, but this approach lacks generality beyond classification within GWAS credible sets. Finally, existing methods rely upon noisy estimates of gene-level association (e.g. from MAGMA or proximity to a GWAS association) to compute gene set enrichments, which degrades sensitivity and specificity of their estimated enrichments and their predicted genes.

Impact of datasets on support calculations

Make sure to add this narrative: that if you don't regularize, you end up with larger gene sets associated. This can actually help prediction in some cases (e.g. gene expression) but harder to interpret. There is a tradeoff and both options should probably be explored


Beyond the need for additional method development, fundamental questions also exist regarding how to most appropriately use gene annotations to prioritize candidate genes. All such methods are biased toward well-studied genes, and the degree to which this bias limits novel discovery is unclear. It is also unclear the extent to which different types of gene annotations (e.g. those derived from animal models vs. those derived from pathway databases vs. those derived from gene expression patterns) lead to different estimates of confidence in a candidate gene, and under which scenarios one type of annotation produces more accurate estimates than another. Many annotations are highly redundant, but the impact of this redundancy on gene set or gene-level analyses has not been well studied – accounting for redundancy might help alleviate false positive results (genes with many redundant annotations) but also, by identifying independent disease-relevant annotations, suggest distinct pathways or disease endophenotypes of high interest in the study of complex diseases such as type 2 diabetes or Alzheimer’s disease.

Embrace circularity -- how does this help us detect and handle it?

PIGEAN can help you *reason*. Can add or remove gene sets. Can say, what would happen if I fix this annotation?

How many independent gene sets are there? For each trait, plot number of non-zero uncorrected vs. number of non-zero beta

Comparison of different assumptions and sources (including lessons for people prioritizing genes)

Run with reconstituted DEPICT sets + the original sets, jointly fit them, see whether the DEPICT sets add anything

Run with just text mining and string naively, show this predicts the GWAS signals but not so much the exomes. Then run with the GWAS removed. Also look outside of GWAS loci. Also look at GWAS published after the text mining
Deploy on textmining based on GWAS -> circular, but still could be useful (summarize what people think)

To find novel genes: (a) run with all, look at only things with high combined but low prior (b) run with all, look at things with high combined but prior only from non mouse/text gene sets (c) run with all, then run with pops, find things high in pops but not in all

How do gene lists change based on different features
Do we observe instances of high in pops but not high in all? Why? Problem with the model or of actually having something else out compete it?

Exomes vs. GWAS in pathways explored. Exomes + GWAS together. Compare results from GWAS vs. results from exomes (same trait), controlling for sample size (inflate GWAS ses to have same number of associations)

All of the annotation sets independently, as well as together (various combinations)

Auto writer of GWAS paper - focus specifically on the gene sets you want (e.g. mouse, MESH). Distinguish application of “what is known” vs. what could be novel.

Single cell prioritization is very popular -- run with just pops, then show beta uncorrected and corrected. Does it do better? Does it do different? When you combine which adjudicates?

Here are all the annotations you think would be good, but actually better captured by others

Here are annotations beware -- string, mesh, because they are circular

Flag genes that have vastly different priors between different collections (what does literature have way different than mice vs. expression/interactomes etc.)

Look at differences S2G and V2G vs. not

LINCS data

Genes with genetic support across 1000 traits

First need to run these:
A2FKP
Rare diseases
Exomes
GCAT top traits


Questions:
Do we add in V2G input. ?
Can we make the reading in more efficient by not using a dictionary?

What genes have most evidence from animals? From literature? From experiments?
What genes do we find from POPS that have no experimental data?
Same thing for other groups of annotations – genes with only one type of support

Look at genes with mouse evidence; what fraction have genetic support? How does it change as function of number? Do this based on beta_uncorrected vs. beta. Highlight genes that actually don’t have support, or those that have more than it would seem
Compare number of mouse annotations (curated) vs. Prior. Identify genes that have actually much less evidence than it would appear (and the converse)
Genes with same sum of beta uncorrected, but different sum of betas. One is validated, one is not
Take all genes with some mouse data -- uncorrected beta above zero. What fraction of those have mouse data after correction? Have combined score? Answers question of all genes with mouse reports, how trustworthy? 
Conversely, take all genes with positive evidence from GWAS and ask how many have mouse data? 

Novel genes, applications of genetic support not limited to effector genes
How many genes were not in GWAS loci? How many genes were in a pathway but not implicated in disease?

Discovery of obvious disease genes even with no significant signals

Architecture ideas:
1. Correlation of traits per vectors of genes (direct and indirect)
2. Same for gene sets (betas and beta uncorrected)
3. UMAP plot
4. Compare to genetic correlations

What fraction of GWAS loci have we solved (one gene with high prior)?

EMPHASIZE RARE/COMMON overlap. Orphanet -> which genes -> common traits; Common -> which genes -> Orphanet

What genes have low huge score but high log bf? These are genes that are novel candidates at a locus and would not have been predicted by closest gene
How many novel genes? How many well known genes? How many WTF genes?

Can we show that genes with high prior but low GWAS are more constrained somehow? Look across all traits for this

Why do some high prior genes not have any signal?
Are large priors but low GWAS enriched for
(a) pli
(B) associations with glucose but not diabetes?

Apply to exome sequencing only
How many additional genes?

Focus on novel effector genes not in the lists we have. Focus on novel associations supported by exomes, or vice versa
Also highlight genes that were significant before but not anymore due to proper adjustment

Drug target support
Drug target prediction (look at Ron Do new article)
Refining the impact of genetic evidence on clinical success
Look up drug targets. How many supported? How many repurposing candidates?
For several diseases, take drugs and year of clinical trial. Then run GWASs up to that using mouse/pubmed data available at that time. See at what year support reached vs. GWAS p-value?
Drug targets analyses -- repeat the recent paper to find our how many have support under this new metric
Known drug targets : see if they have support from the single cell data. Reference decent paper and debate on twitter

Monogenic disease support

Indirect support: genes that are in pathways for a disease.
Without PIGEAN, you would look for enriched pathways, and then check overlap. But how do you combine multiple pathways for a gene?

Differences between genes with and without direct support. Are they fundamentally different in any way? Constraint? What is known about them?

For genes with no indirect support but direct support, looking at other traits for which they do have indirect support could be very informative about new mechanisms for the disease

Analysis of the "dark GWAS":
- How many GWAS loci have a gene with a known disease mechanism?
- How many GWAS loci have a gene with a known mechanism for a correlated trait?
- How many GWAS loci have a gene with a known mechanism for any trait?
- How many are completely novel?


For mono and drug, look at pathway support vs direct support. Mono might have more pathway support than direct support

Example query: upload list of DE genes or islet screen genes, calculate support

Gene correlations based on combined, log_bf, or prior. Which genes have similar partners change?
Which gene is most pleiotropic? Overall and across PCs

Distribution of number of traits across genes/gene sets (after independent correction) -- is it bimodal or continuous?

How many / what fraction of genes have known mechanisms? Does this vary across traits?

Application: for all of the GWAS with effector gene lists, how accurate were we?

PPARG maternal birth weight. Just below sig, becomes high. Even better: PPARG fetal birth weight. No signal at all but prior is super high
CD19 IBD
CAD and lipids vs. vascular
WHR vs. WHRadjBMI – the adjBMI doesn’t have the metabolic clusters

Tools for filtering novel vs known associations.
Annotating GWAS produced by a methods paper (when you don’t have expertise in the trait you are analyzing)

Reasons for trial stoppage


Indirect support (constraint?) or broad.expression leads to safety concerns


Efficacy due to lack of support


Should we redo these with pigean?


https://www.nature.com/articles/s41588-024-01854-z

Use of indirect support for target validation

Redo Minikel

Even if it is biased, it is a scaffold to understand mechanisms
Many genes in the pathways don't have indirect support
(Show example of genes in gene sets with high beta uncorrected but not high beta)


Filling in missing data

Of the pathways we identify, how many signals are explained vs. novel?


How many processes are tagged by mouse phenotypes?


How many mouse phenotypes should be screened for each trait?


That is, what experiments "complete the picture"


—-

Methods

PIGEAN model

HuGE score extensions and calculations

Inferring parameters from data
How we do p
How we do alpha
Sigma: doesn’t work when try to estimate like LD-pred, so we tried just max beta. This works okay but is sensitive so we instead empirically chose something based upon runs across many GWAS

Tuning regularization parameters
- sigma
- sigma_power 
- for default model (pow0), a sigma of 2 less than simulated increases accuracy, although doesn’t quite match size
- it does match size for powc model, although it decreases accuracy of genes slightly
- since we don’t know true model, we think 2 is a good default since it matches size and/or increases accuracy for powers between 0 and c

Interpretation of PIGEAN outputs

Simulating data
- truncated model, null GWAS, non-null GWAS

Permuted data


\end{comment}

\end{document}
